\documentclass[11pt]{amsart}

\usepackage[margin=1.1in]{geometry}
\usepackage{amsmath,amssymb,amsthm,mathtools}
\usepackage{enumitem}
\usepackage{microtype}
\usepackage{mathrsfs}
\usepackage{hyperref}
\usepackage[nameinlink,capitalize]{cleveref}

\hypersetup{
  colorlinks=true,
  linkcolor=blue,
  citecolor=blue,
  urlcolor=blue
}

\numberwithin{equation}{section}

\theoremstyle{plain}
\newtheorem{theorem}{Theorem}[section]
\newtheorem{proposition}[theorem]{Proposition}
\newtheorem{lemma}[theorem]{Lemma}
\newtheorem{corollary}[theorem]{Corollary}

\theoremstyle{definition}
\newtheorem{definition}[theorem]{Definition}

\theoremstyle{remark}
\newtheorem{remark}[theorem]{Remark}

\newcommand{\R}{\mathbb R}
\newcommand{\N}{\mathbb N}
\newcommand{\eps}{\varepsilon}
\newcommand{\loc}{\mathrm{loc}}
\newcommand{\CKN}{\mathrm{CKN}}
\newcommand{\NS}{\mathrm{NS}}
\newcommand{\RNSE}{\mathrm{RNSE}}
\newcommand{\divergence}{\operatorname{div}}
\newcommand{\dist}{\operatorname{dist}}
\newcommand{\supp}{\operatorname{supp}}
\newcommand{\esssup}{\operatorname*{ess\,sup}}
\newcommand{\calS}{\mathcal S}
\newcommand{\calR}{\mathcal R}
\newcommand{\calX}{\mathcal X}
\newcommand{\calU}{\mathcal U}
\newcommand{\calE}{\mathcal E}
\newcommand{\calA}{\mathcal A}
\newcommand{\calM}{\mathcal M}
\newcommand{\calF}{\mathcal F}
\newcommand{\calC}{\mathcal C}
\newcommand{\calG}{\mathcal G}
\newcommand{\calK}{\mathcal K}
\newcommand{\Cax}{\mathcal C_{\mathrm{ax}}}
\newcommand{\Crot}{\mathcal C_{\mathrm{rot}}}
\newcommand{\Cstat}{\mathcal C_{\mathrm{stat}}^{\mathrm{ret}}}
\newcommand{\Cgen}{\mathcal C_{\mathrm{gen}}}
\newcommand{\bP}{\mathbb P}

\title[Residual Type I Seregin limit exclusion]
{Local Exclusion of Residual Type I Navier--Stokes Seregin Limits}

\author{Guillem Duran Ballester}
\date{}

\subjclass[2020]{35Q30, 35B44, 35B40, 76D05}
\keywords{Navier--Stokes equations, Type I singularities, ancient solutions, Seregin limits, Liouville theorems, local concentration}

\begin{document}

\begin{abstract}
We assemble the refined residual analysis for Type I
Navier--Stokes Seregin limits.  A finite-energy Type I singularity produces,
through the Seregin blow-up procedure, a nonzero bounded ancient solution of
the centered self-similar Navier--Stokes equation carrying positive local
velocity \(L^3\)-mass.  After the
previously treated small, stationary \(L^3\), uniformly \(L^3\)-tight, and
closed structured/decay alternatives have been removed, the remaining
residual class is divided into four analytic classes: axisymmetric
bounded-circulation profiles, rotational relative equilibria, degenerate
stationary-hull profiles with retained local concentration, and the generic
terminal residual.

The first class is excluded by a centered angular-circulation equation and an
axis-absorption Liouville theorem.  The second class is closed by a local
Coriolis-flux identity and a local concentration decomposition: a superthreshold
rotational flux gives a nontrivial recentered local profile, while a subthreshold flux
remains inside the compact local limiting equation.  The third class is
reduced, by scale reselection for time-translation hulls, to a stationary
profile in a degenerate stationary family which occurs as a Seregin limit
and still carries local CKN concentration; it is then excluded by the terminal
local concentration theorem.  The final class is
reduced by covariant tail normalization, exclusion of finite separated
concentration profiles, indecomposable-profile sequence-\(L^3\) completeness,
mildness inheritance, and the Albritton--Barker endpoint Liouville theorem.  The paper
distinguishes statements imported from Papers I--III and statements proved
here.  No global Coriolis--Leray Liouville theorem for bounded profiles or global
tail estimate is used in the residual closure.
\end{abstract}

\maketitle
\tableofcontents

\section{Introduction and main theorem}
\label{sec:introduction}

We write the incompressible Navier--Stokes equations on
\(\R^3\times(0,T)\) as
\begin{equation}\label{eq:NS-physical}
   \partial_tu-\Delta u+(u\cdot\nabla)u+\nabla p=0,
   \qquad \nabla\cdot u=0 .
\end{equation}
The residual analysis in this paper begins after the standard Type I blow-up
construction has been performed.  Thus the basic solution is a smooth bounded
ancient solution of the centered equation
\begin{equation}\label{eq:centered-main}
   \partial_\tau V-\Delta V+\frac12 y\cdot\nabla V+\frac12 V
   +(V\cdot\nabla)V+\nabla P=0,
   \qquad \nabla\cdot V=0 ,
\end{equation}
on \(\R^3_y\times\R_\tau\).  The equation is obtained from the physical
ancient representative \((U,Q)\) by
\begin{equation}\label{eq:centered-transform}
   y=\frac{x}{\sqrt{-t}},\qquad
   \tau=-\log(-t),\qquad
   V(y,\tau)=\sqrt{-t}\,U(x,t),
   \qquad
   P(y,\tau)=(-t)Q(x,t),
\end{equation}
where \(t<0\) and pressures are always understood modulo addition of a
function of time.  Conversely,
\begin{equation}\label{eq:physical-pullback-main}
   U(x,t)=(-t)^{-1/2}
   V\!\left(\frac{x}{\sqrt{-t}},-\log(-t)\right),
   \qquad t<0 .
\end{equation}

The purpose of the paper is to present the refined residual decomposition as a
single standard PDE argument, using only the reductions already established in
the preceding papers of the series.  To keep the logical status clear, we first state the imported Paper
0 and Paper I--III dependencies.  Apart from those four companion papers, the
only external inputs used below are the cited peer-reviewed Navier--Stokes
regularity and Liouville theorems.  All rotational-flux estimates, terminal
local concentration decompositions, exclusions of local compactness failure and
diffuse exterior concentration, and finite separated-profile exclusions used
below are proved in the present paper from local estimates rather than imported
as global estimates.

\begin{theorem}[Paper I Type I Seregin extraction theorem]
\label{hyp:base-seregin-package}
A finite-energy Type I singular point produces a normalized Seregin ancient
limit \((V,P)\) satisfying \eqref{eq:centered-main}, with
\[
   \|V\|_{L^\infty(\R^3\times\R)}\le M<\infty,
\]
local suitability and local pressure-gauge compactness, and a compact
parabolic cylinder \(K\Subset\R^3\times\R\) for which
\begin{equation}\label{eq:positive-velocity-main}
   \iint_K |V|^3\,dy\,d\tau
   \ge \eta_0>0 .
\end{equation}
The collection of all normalized Seregin limits generated by the original
solution is
denoted by \(\calS\).
\end{theorem}

\begin{proof}
This is Paper I's extraction theorem \cite{DuranPaperI}, including velocity
nonvanishing, local suitability, and local pressure compactness.  The underlying
positive local concentration input comes from Paper 0 \cite{DuranPaper0}.  The
present paper uses this as an imported reduction and does not reprove the
Type I blow-up compactness theorem.
\end{proof}

\begin{definition}[Imported earlier alternatives]
\label{hyp:previous-nonresidual}
Within a fixed normalized Seregin collection \(\calS\), the small-amplitude,
stationary \(L^3\), uniformly \(L^3\)-tight, and closed structured/decay
classes are the non-residual alternatives defined in Paper I.  The
small-amplitude and stationary \(L^3\) alternatives are closed in Paper I; the
uniformly tight class is excluded in Paper II \cite{DuranPaperII}; and the
closed structured/decay class is precisely the part of the structured
alternatives proved empty in Paper III \cite{DuranPaperIII}.  Any
structured profile not belonging to that closed Paper I class is residual by
definition and is handled by the present paper.  The present paper does not
reprove those earlier alternatives and does not add any hypothesis to their
statements.
\end{definition}

\begin{definition}[Coarse residual class]
\label{def:coarse-residual}
Let \(I,J\) be the index sets used in the base Type I residual
criterion.  The coarse residual class
\(\calR(\calS;I,J)\) is the complement in \(\calS\) of the already removed
alternatives in \Cref{hyp:previous-nonresidual}.
\end{definition}

\begin{definition}[Refined residual decomposition]
\label{def:refined-decomposition}
The refined residual decomposition is the ordered decomposition of the remaining
coarse residual class into the following alternatives, each understood after
all earlier alternatives and all previously closed non-residual classes have
been removed:
\[
   \Cax,\qquad
   \Crot,\qquad
   \Cstat,\qquad
   \Cgen .
\]
Here \(\Cax\) is the axisymmetric bounded-circulation class, \(\Crot\) is the
rotational relative-equilibrium class, \(\Cstat\) is the degenerate
stationary-hull class with retained local concentration, and \(\Cgen\) is the
generic terminal residual class.
\end{definition}

\begin{theorem}[Refined residual closure by local concentration analysis]
\label{thm:refined-residual-closure}
Relative to the imported Paper I--III alternatives in
\Cref{hyp:base-seregin-package,hyp:previous-nonresidual}, one has
\[
   \Cax=\varnothing,\qquad
   \Crot=\varnothing,\qquad
   \Cstat=\varnothing,\qquad
   \Cgen=\varnothing .
\]
Consequently the coarse residual class \(\calR(\calS;I,J)\) is empty.
\end{theorem}

\begin{proof}
The axisymmetric bounded-circulation class is excluded by
\Cref{thm:R1-exclusion}.  The rotational relative-equilibrium class is excluded by
\Cref{thm:R2-exclusion-safe}, which uses the localized Coriolis-flux identity
and a terminal local concentration decomposition rather than a whole-space
Coriolis--Leray Liouville theorem.
The degenerate stationary-hull class with retained local concentration
\(\Cstat\) is excluded by \Cref{thm:R3S-R3-empty}, using the
terminal local concentration theorem proved in
\Cref{sec:terminal-concentration,sec:separated-profiles-atomic}.
Finally, the generic terminal residual class \(\Cgen\) is excluded by
\Cref{thm:R4-final-closure}.  These four ordered exclusions exhaust the
refined residual decomposition.
\end{proof}

\begin{corollary}[Paper I residual input]
\label{cor:paperI-residual-input}
For every normalized Seregin collection \(\calS\) generated by a
Paper~0-admissible finite-energy Type I singularity, the Paper I remaining
class is empty.
\end{corollary}

\begin{proof}
Paper I defines its remaining class as the complement, inside \(\calS\), of
the small-amplitude, stationary \(L^3\), uniformly \(L^3\)-tight, and closed
structured/decay classes.  The notation \(\calR(\calS;I,J)\) used here records
the same complement together with the ordered removals used in the residual
decomposition; the indices \(I,J\) add no analytic assumption.  By
\Cref{thm:refined-residual-closure}, this class is empty.
\end{proof}

\begin{theorem}[Type I exclusion after local residual closure]
\label{thm:typeI-exclusion-main}
No Paper~0-admissible finite-energy Type I singularity exists whose normalized
Seregin collection is \(\calS\).
\end{theorem}

\begin{proof}
Suppose that a Type I singular point exists.  By
\Cref{hyp:base-seregin-package}, the Seregin blow-up procedure produces a
normalized ancient limit \((V,P)\) satisfying \eqref{eq:centered-main} and the
local velocity nonvanishing condition \eqref{eq:positive-velocity-main}.  By
the Paper I class exhaustion, \(V\) is either in one of the non-residual
classes recorded in \Cref{hyp:previous-nonresidual} or in
\(\calR(\calS;I,J)\).  The first alternative is impossible by the corresponding
closed inputs: the small and stationary \(L^3\) classes are excluded in Paper I,
the uniformly \(L^3\)-tight class is excluded in Paper II, and the closed
structured/decay class is excluded in Paper III.  The second alternative is
impossible by \Cref{thm:refined-residual-closure}.  This contradiction proves
the theorem.
\end{proof}

\section{Shared local notation}
\label{sec:shared-notation}

For \(z_0=(y_0,\tau_0)\) and \(r>0\), write
\[
   Q_r(z_0):=B_r(y_0)\times(\tau_0-r^2,\tau_0).
\]
For a centered profile \((V,P)\), the pressure-normalized local critical
quantity is
\begin{equation}\label{eq:local-energy-main}
   \calE_r(V,P;z_0)
   :=
   \iint_{Q_r(z_0)}
   \left(|V|^3+|P-c_{r,z_0}(\tau)|^{3/2}\right)\,dy\,d\tau,
\end{equation}
where \(c_{r,z_0}(\tau)\) is an allowed time-dependent pressure gauge, usually
chosen as a spatial average over \(B_r(y_0)\).  All local pressure
convergences below are understood after such gauges.

\begin{definition}[Retained local concentration profile]
\label{def:retained-active-main}
Fix \(\eps_*>0\).  A terminal profile \((U,\Pi)\) has retained local
concentration if, after the terminal normalization under consideration,
\[
   \calE_1(U,\Pi;0)\ge \eps_* .
\]
\end{definition}

\begin{definition}[Refined residual candidate]
\label{def:refined-residual-candidate}
A refined residual candidate is a bounded normalized Seregin ancient limit
obtained from the Paper I extraction theorem which carries the positive compact local
velocity mass in \eqref{eq:positive-velocity-main} and which is not in any alternative
already closed before the refined residual decomposition is entered.  The definition is local:
the only persistent lower bound attached to the candidate by Paper I is the
fixed compact velocity \(L^3\) lower bound.
\end{definition}

\begin{remark}[No hidden global \(L^3\) assertion]
\label{rem:no-hidden-global-L3}
The terminal arguments below do not claim that a generic residual profile is
globally \(L^3\)-tight or globally bounded in \(L^\infty_\tau L^3_y\).  The
only nonlocal critical-norm conclusion used at the endpoint is the weaker
sequence-\(L^3\) completeness of retained indecomposable concentration
profiles, and that conclusion is derived from local terminal alternatives rather than from an
a priori global estimate on the original residual class.
\end{remark}

\section{Exclusion of the axisymmetric bounded-circulation residual class}
\label{sec:R1-exclusion}

In this section we prove the rigidity statement needed for the first refined
residual class.  The argument uses the scale-invariant angular circulation in
centered variables.  The essential point is that the centered Type~I equation
contains the Ornstein--Uhlenbeck drift.  Thus the angular-circulation equation
has an absorbing boundary at the symmetry axis together with a confining drift
at spatial infinity.  This gives a scalar Liouville theorem for bounded ancient
solutions of the circulation equation.

Throughout the section we write cylindrical coordinates as
\((\rho,\theta,Z)\), and an axisymmetric vector field as
\[
   V=V_\rho e_\rho+V_\theta e_\theta+V_Z e_Z .
\]
The centered Navier--Stokes equation is
\begin{equation}\label{eq:centered-R1}
   \partial_\tau V-\Delta V+\frac12 y\cdot\nabla V+\frac12 V
   +(V\cdot\nabla)V+\nabla P=0,
   \qquad \nabla\cdot V=0 .
\end{equation}
For an axisymmetric solution define the centered angular circulation
\begin{equation}\label{eq:centered-Gamma-def}
   \Gamma(\rho,Z,\tau):=\rho V_\theta(\rho,Z,\tau).
\end{equation}
This is the same scale-invariant quantity as the physical circulation
\(rU_\theta\).  If the physical representative is
\[
   U(x,t)=(-t)^{-1/2}V(x/\sqrt{-t},-\log(-t)),
\]
then
\[
   rU_\theta(r,z,t)=\rho V_\theta(\rho,Z,\tau),
   \qquad r=\sqrt{-t}\,\rho,
   \quad z=\sqrt{-t}\,Z.
\]

\begin{lemma}[Centered circulation equation]\label{lem:centered-circulation-equation}
Let \((V,P)\) be a smooth bounded axisymmetric solution of
\eqref{eq:centered-R1}.  Then \(\Gamma=\rho V_\theta\) satisfies
\begin{equation}\label{eq:centered-Gamma-equation}
   \partial_\tau\Gamma
   =
   \partial_\rho^2\Gamma-\frac1\rho\partial_\rho\Gamma
      +\partial_Z^2\Gamma
   -\left(V_\rho+\frac\rho2\right)\partial_\rho\Gamma
   -\left(V_Z+\frac Z2\right)\partial_Z\Gamma
\end{equation}
on \((0,\infty)\times\mathbb R\times\mathbb R\).  Moreover
\begin{equation}\label{eq:axis-boundary-Gamma}
   \Gamma(0,Z,\tau)=0
   \qquad\hbox{for all }(Z,\tau)\in\mathbb R\times\mathbb R .
\end{equation}
\end{lemma}

\begin{proof}
The angular component of \eqref{eq:centered-R1} is
\[
   \partial_\tau V_\theta
   -\left(\Delta-\frac1{\rho^2}\right)V_\theta
   +\frac12(\rho\partial_\rho+Z\partial_Z)V_\theta
   +\frac12V_\theta
   +V_\rho\partial_\rho V_\theta+V_Z\partial_ZV_\theta
   +\frac{V_\rho V_\theta}{\rho}=0 .
\]
Multiplying by \(\rho\), using
\[
   \rho\left(\Delta-\frac1{\rho^2}\right)V_\theta
   =\partial_\rho^2\Gamma-\frac1\rho\partial_\rho\Gamma+
      \partial_Z^2\Gamma,
\]
\[
   \rho\left[\frac12(\rho\partial_\rho+Z\partial_Z)V_\theta
      +\frac12V_\theta\right]
   =\frac12(\rho\partial_\rho+Z\partial_Z)\Gamma,
\]
and
\[
   \rho\left(V_\rho\partial_\rho V_\theta+V_Z\partial_ZV_\theta
      +\frac{V_\rho V_\theta}{\rho}\right)
   =V_\rho\partial_\rho\Gamma+V_Z\partial_Z\Gamma,
\]
gives \eqref{eq:centered-Gamma-equation}.  Finally, smoothness of a bounded
axisymmetric vector field at the axis implies \(V_\theta(\rho,Z,\tau)=O(\rho)\)
as \(\rho\downarrow0\), locally uniformly in \((Z,\tau)\).  Hence
\(\Gamma=\rho V_\theta=O(\rho^2)\), which gives
\eqref{eq:axis-boundary-Gamma}.
\end{proof}

We next record the scalar Liouville theorem used below.  It is stated in a
form that isolates precisely the PDE mechanism: bounded drift perturbations of
the centered circulation operator cannot sustain a bounded ancient solution
which vanishes on the symmetry axis.

\begin{lemma}[Axis-absorption Liouville theorem]
\label{lem:axis-absorption-liouville}
Let \(B=(B_\rho,B_Z)\) be a smooth vector field on
\((0,\infty)\times\mathbb R\times\mathbb R\) satisfying
\[
   \|B\|_{L^\infty}<\infty .
\]
Let \(G\in C^{2,1}((0,\infty)\times\mathbb R\times\mathbb R)\cap L^\infty\)
solve
\begin{equation}\label{eq:abstract-Gamma-equation}
   \partial_\tau G
   =
   \partial_\rho^2G-\frac1\rho\partial_\rho G+
      \partial_Z^2G
   -\left(B_\rho+\frac\rho2\right)\partial_\rho G
   -\left(B_Z+\frac Z2\right)\partial_ZG
\end{equation}
for all \(\tau\in\mathbb R\), and assume
\begin{equation}\label{eq:abstract-axis-zero}
   \lim_{\rho\downarrow0}G(\rho,Z,\tau)=0
   \qquad\hbox{locally uniformly in }(Z,\tau).
\end{equation}
Then \(G\equiv0\).
\end{lemma}

\begin{proof}
Set \(M_B=\|B\|_{L^\infty}\).  For \(R>1\) define
\begin{equation}\label{eq:radial-barrier}
   h_R(\rho)
   :=
   \frac{\displaystyle\int_0^{\min\{\rho,R\}} s
      \exp\left(\frac{s^2}{4}-M_Bs\right)\,ds}
   {\displaystyle\int_0^R s
      \exp\left(\frac{s^2}{4}-M_Bs\right)\,ds} .
\end{equation}
Then \(0\le h_R\le1\), \(h_R(0)=0\), \(h_R(R)=1\), and, on
\(0<\rho<R\),
\begin{equation}\label{eq:hR-ode}
   h_R''+\left(M_B-\frac\rho2-\frac1\rho\right)h_R'=0,
   \qquad h_R'>0 .
\end{equation}
Let
\[
   \mathcal L_\tau f
   :=
   \partial_\rho^2f-\frac1\rho\partial_\rho f+\partial_Z^2f
   -\left(B_\rho+\frac\rho2\right)\partial_\rho f
   -\left(B_Z+\frac Z2\right)\partial_Z f .
\]
Since \(B_\rho\ge -M_B\), \eqref{eq:hR-ode} gives
\begin{equation}\label{eq:hR-superharmonic}
   \mathcal L_\tau h_R
   =h_R''-\left(\frac1\rho+\frac\rho2+B_\rho\right)h_R'
   \le
   h_R''+\left(M_B-\frac\rho2-\frac1\rho\right)h_R'
   =0 .
\end{equation}
Thus \((\partial_\tau-\mathcal L_\tau)h_R\ge0\).

We now compare \(G\) to this barrier on the strip
\(\Omega_R=(0,R)\times\mathbb R\).  Let
\(M_G=\|G\|_{L^\infty}\).  We first record the killed-strip decay estimate used
by the comparison argument.  If \(q_{s,R}\) solves
\[
   \partial_\tau q_{s,R}=\mathcal L_\tau q_{s,R},
   \qquad q_{s,R}=0\hbox{ on }\rho=0,R,
   \qquad q_{s,R}(\rho,Z,s)=1,
\]
then, for every compact \(K\Subset\Omega_R\) and every fixed \(t\),
\begin{equation}\label{eq:killed-strip-decay}
   \lim_{s\to-\infty}\sup_{(\rho,Z)\in K}q_{s,R}(\rho,Z,t)=0 .
\end{equation}
To prove \eqref{eq:killed-strip-decay}, fix \(K\Subset\Omega_R\).  Choose
\(\delta>0\), \(Z_K>0\), and \(Z_1>Z_K+1\) so that
\[
   K\subset [\delta,R-\delta]\times[-Z_K,Z_K]
   \subset (0,R)\times(-Z_1,Z_1).
\]
On the finite cylinder
\[
   \Omega_{R,Z_1}:=(0,R)\times(-Z_1,Z_1)
\]
with Dirichlet data on the full parabolic boundary, the operator
\(\mathcal L_\tau\) is uniformly parabolic on
\([\delta/2,R-\delta/2]\times[-Z_1,Z_1]\), has bounded first-order coefficients
there, and has absorbing boundary at \(\rho=0,R\).  The boundary parabolic
Harnack inequality and the strong maximum principle give a time \(T_R>0\) and a
number \(\vartheta_R\in(0,1)\), depending only on
\(R,M_B,\delta,Z_K\) and not on \(Z_1\), such that every nonnegative solution
with initial data bounded by \(1\) and zero lateral data satisfies
\[
   0\le q(\cdot,s+T_R)\le 1-\vartheta_R
   \qquad\hbox{on }[\delta,R-\delta]\times[-Z_K,Z_K].
\]
Applying the same estimate after each time translation and using the maximum
principle between the compact set and the boundary gives
\[
   \sup_K q_{s,R}(\cdot,s+mT_R)
   \le (1-\vartheta_R)^m,\qquad m=1,2,\dots .
\]
For a fixed \(t>s\), choose \(m\) with
\(s+mT_R\le t<s+(m+1)T_R\).  The maximum principle and interior Harnack estimate
on the final interval \([s+mT_R,t]\) preserve the preceding bound up to a
constant depending only on \(R,M_B,\delta,Z_K,T_R\).  Letting \(s\to-\infty\)
forces \(m\to\infty\), proving decay on \(K\) for the truncated problem.  The
solutions on \(\Omega_{R,Z_1}\) increase monotonically to the killed-strip
solution as \(Z_1\to\infty\); passing to this limit gives
\eqref{eq:killed-strip-decay}.

Fix \(R>1\), \(s<t\), and set
\[
   W^+_{R}:=\frac{G}{M_G}-h_R,
   \qquad
   W^-_{R}:=-\frac{G}{M_G}-h_R,
\]
with the convention that the conclusion is immediate if \(M_G=0\).
By \eqref{eq:abstract-Gamma-equation} and \eqref{eq:hR-superharmonic},
\[
   (\partial_\tau-\mathcal L_\tau)W_R^\pm\le0
   \qquad\hbox{in }\Omega_R\times(s,t).
\]
On \(\rho=0\), \eqref{eq:abstract-axis-zero} and \(h_R(0)=0\) give
\(W_R^\pm\le0\).  On \(\rho=R\), \(h_R(R)=1\) and \(|G|\le M_G\) give again
\(W_R^\pm\le0\).  The parabolic comparison principle with the killed-strip
kernel therefore yields, for \(\tau=t\),
\[
   (W_R^\pm)_+(\rho,Z,t)
   \le q_{s,R}(\rho,Z,t)
   \qquad ((\rho,Z)\in\Omega_R).
\]
Letting \(s\to-\infty\) and using \eqref{eq:killed-strip-decay} gives
\(W_R^\pm\le0\) on compact subsets of \(\Omega_R\).  Since the compact subset
is arbitrary,
\begin{equation}\label{eq:G-bounded-by-hR}
   |G(\rho,Z,t)|\le M_G h_R(\rho)
   \qquad (0<\rho<R, Z\in\mathbb R,
   \ t\in\mathbb R).
\end{equation}
Finally, for each fixed \(\rho<\infty\), the denominator in
\eqref{eq:radial-barrier} tends to infinity super-exponentially as
\(R\to\infty\), while the numerator remains fixed.  Hence
\(h_R(\rho)\to0\).  Letting \(R\to\infty\) in
\eqref{eq:G-bounded-by-hR} gives \(G(\rho,Z,t)=0\) for every
\(\rho>0\), \(Z\in\mathbb R\), and \(t\in\mathbb R\).  The boundary value at
\(\rho=0\) is already zero by assumption.  Therefore \(G\equiv0\).
\end{proof}

\begin{theorem}[Exclusion of the axisymmetric bounded-circulation residual class]
\label{thm:R1-exclusion}
Let \(\Cax(\mathcal S;I,J)\) be the axisymmetric bounded-circulation residual
class in the refined residual decomposition.  Then
\begin{equation}\label{eq:R1-empty}
   \Cax(\mathcal S;I,J)=\varnothing .
\end{equation}
\end{theorem}

\begin{proof}
Assume, for contradiction, that
\(V\in\Cax(\mathcal S;I,J)\).  By definition of \(\Cax\), the
associated physical ancient representative is axisymmetric and has bounded
angular circulation.  Equivalently, in centered variables,
\[
   \Gamma(\rho,Z,\tau)=\rho V_\theta(\rho,Z,\tau)
\]
is bounded on \((0,\infty)\times\mathbb R\times\mathbb R\).  Since \(V\) is a
Seregin ancient limit, it is smooth and bounded in centered variables.  Hence
\(B=(V_\rho,V_Z)\) is bounded.

By \Cref{lem:centered-circulation-equation}, \(\Gamma\) satisfies
\eqref{eq:centered-Gamma-equation} and vanishes at the axis.  Applying
\Cref{lem:axis-absorption-liouville} with \(G=\Gamma\) and
\(B=(V_\rho,V_Z)\) gives
\[
   \Gamma\equiv0 .
\]
Therefore \(V_\theta\equiv0\) for \(\rho>0\), and by smoothness also on the
axis.  Thus the associated ancient solution is axisymmetric without swirl.
But the axisymmetric no-swirl class is one of the structured alternatives
removed before the residual class is defined.  This contradicts
\(V\in\mathcal R(\mathcal S;I,J)\), and therefore no such \(V\) exists.
\end{proof}

\begin{remark}[Why this does not assert the full bounded-swirl Liouville theorem]
The proof uses the centered Type~I equation, specifically the
Ornstein--Uhlenbeck drift terms \(-\rho\partial_\rho/2\) and
\(-Z\partial_Z/2\) in the scalar equation for \(\Gamma\).  It is therefore a
Type~I Seregin limit argument.  It does not prove that every bounded ancient
axisymmetric solution of the unrenormalized Navier--Stokes equations with
bounded physical circulation is swirl-free.  The latter lacks the confining
centered drift and remains a different shell-transport problem.
\end{remark}


\section{The rotational relative-equilibrium residual class}
\label{sec:R2-reduction}

We next isolate the second refined residual class.  This class consists of
centered Type~I ancient limits whose time dependence is a rigid rotation of a
fixed spatial profile.  Passing to the co-rotating frame converts the centered
Navier--Stokes equation into a stationary Leray profile equation with a
Coriolis transport term.  A whole-space Coriolis--Leray Liouville theorem for
all bounded profiles is not used here.  Instead, the Coriolis term is localized as an annular
flux and is then treated by the local concentration analysis proved in
\Cref{sec:terminal-concentration,sec:separated-profiles-atomic}.

Throughout this section the centered equation is
\begin{equation}\label{eq:centered-R2-safe}
   \partial_\tau V-\Delta V+\frac12 y\cdot\nabla V+\frac12 V
   +(V\cdot\nabla)V+\nabla P=0,
   \qquad \nabla\cdot V=0 .
\end{equation}
For \(\Omega\in\mathfrak{so}(3)\), set
\begin{equation}\label{eq:rotation-group-R2-safe}
   Q(\tau):=\exp(\tau\Omega)\in SO(3).
\end{equation}

\begin{definition}[Rotational relative-equilibrium profile]
\label{def:R2-relative-profile-safe}
A bounded smooth ancient solution \((V,P)\) of \eqref{eq:centered-R2-safe} is a
rotational relative equilibrium if there exist
\(\Omega\in\mathfrak{so}(3)\), a bounded smooth divergence-free vector field
\(W\), and a pressure profile \(\Pi\), such that
\begin{equation}\label{eq:R2-relative-form-safe}
   V(y,\tau)=Q(\tau)W(Q(\tau)^Ty),
   \qquad
   P(y,\tau)=\Pi(Q(\tau)^Ty)
\end{equation}
up to the usual time-dependent pressure gauge.
The class \(\Crot\) consists of such residual Seregin profiles, after the
axisymmetric bounded-circulation class and the previously closed Paper I--III
alternatives have been removed.
\end{definition}

\begin{lemma}[Co-rotating stationary equation]
\label{lem:co-rotating-equation-safe}
Let \((V,P)\) be a rotational relative equilibrium.  Then the co-rotating
profile \((W,\Pi)\) satisfies
\begin{equation}\label{eq:R2-coriolis-profile-safe}
   -\Delta W+\frac12 y\cdot\nabla W+\frac12 W
   +(W\cdot\nabla)W
   +\Omega W-(\Omega y)\cdot\nabla W
   +\nabla\Pi=0,
   \qquad \nabla\cdot W=0 .
\end{equation}
Equivalently,
\begin{equation}\label{eq:R2-coriolis-profile-project-safe}
   \mathbb P\Big[
   -\Delta W+\frac12 y\cdot\nabla W+\frac12 W
   +(W\cdot\nabla)W
   +\Omega W-(\Omega y)\cdot\nabla W
   \Big]=0,
\end{equation}
where \(\mathbb P\) is the whole-space Leray projector.
\end{lemma}

\begin{proof}
Let \(z=Q(\tau)^Ty\).  Since \(Q'(\tau)=\Omega Q(\tau)\) and
\(Q(\tau)^T\Omega Q(\tau)=\Omega\), one has
\[
   \partial_\tau z=-\Omega z.
\]
Therefore
\[
   \partial_\tau V(y,\tau)
   =Q(\tau)\Big(\Omega W(z)-(\Omega z)\cdot\nabla W(z)\Big).
\]
The Laplacian, divergence, convection term, and pressure gradient are invariant
under the orthogonal change of variables.  Substitution into
\eqref{eq:centered-R2-safe}, followed by multiplication by \(Q(\tau)^T\), gives
\eqref{eq:R2-coriolis-profile-safe}.  Applying the Leray projector gives
\eqref{eq:R2-coriolis-profile-project-safe}.
\end{proof}

\subsection{The Bernoulli identity and the local Coriolis flux}

Write
\begin{equation}\label{eq:R2-F-def}
   F(y):=\frac12y-\Omega y,
   \qquad
   A(y):=W(y)+F(y),
   \qquad
   \omega:=\nabla\times W .
\end{equation}
Since \(\Omega\) is skew-symmetric, there is a unique vector
\(\varpi\in\mathbb R^3\) such that
\begin{equation}\label{eq:R2-Omega-vector}
   \Omega y=\varpi\times y .
\end{equation}
Then
\begin{equation}\label{eq:R2-curlF}
   \nabla\times F=-2\varpi .
\end{equation}

\begin{lemma}[Coriolis Bernoulli identity]
\label{lem:R2-coriolis-bernoulli}
Let \((W,\Pi)\) be a smooth solution of
\eqref{eq:R2-coriolis-profile-safe}.  Define
\begin{equation}\label{eq:R2-Bernoulli-def}
   B:=\frac12|W|^2+\Pi+F\cdot W .
\end{equation}
Then
\begin{equation}\label{eq:R2-gradient-B}
   \nabla B+\nabla\times\omega-A\times\omega=0,
\end{equation}
and consequently
\begin{equation}\label{eq:R2-Bernoulli-PDE}
   -\Delta B+A\cdot\nabla B
   =
   -|\omega|^2-\omega\cdot(\nabla\times F)
   =
   -|\omega|^2+2\varpi\cdot\omega .
\end{equation}
\end{lemma}

\begin{proof}
The vector identity
\[
   (W\cdot\nabla)W=\nabla\left(\frac12|W|^2\right)-W\times\omega
\]
and the identity
\[
   (F\cdot\nabla)W+(\nabla F)^TW
   =\nabla(F\cdot W)-F\times\omega
\]
apply because
\((\nabla F)^TW=\frac12W+\Omega W\).  Thus
\eqref{eq:R2-coriolis-profile-safe} can be rewritten as
\[
   -\Delta W+
   \nabla\left(\frac12|W|^2+\Pi+F\cdot W\right)
   -(W+F)\times\omega=0 .
\]
Since \(\nabla\cdot W=0\), one has
\(-\Delta W=\nabla\times\omega\), which gives
\eqref{eq:R2-gradient-B}.

Taking divergence in \eqref{eq:R2-gradient-B} gives
\[
   \Delta B=\nabla\cdot(A\times\omega),
\]
because the divergence of a curl vanishes.  Using
\[
   \nabla\cdot(A\times\omega)
   =\omega\cdot(\nabla\times A)-A\cdot(\nabla\times\omega),
   \qquad
   \nabla\times A=\omega+\nabla\times F,
\]
we obtain
\[
   \Delta B
   =|\omega|^2+\omega\cdot(\nabla\times F)
      -A\cdot(\nabla\times\omega).
\]
On the other hand, by \eqref{eq:R2-gradient-B},
\[
   A\cdot\nabla B
   =A\cdot(A\times\omega)-A\cdot(\nabla\times\omega)
   =-A\cdot(\nabla\times\omega).
\]
Combining the last two identities gives \eqref{eq:R2-Bernoulli-PDE}.  Finally,
\eqref{eq:R2-curlF} follows from \eqref{eq:R2-Omega-vector}.
\end{proof}

\begin{lemma}[Localized Coriolis flux identity]
\label{lem:R2-coriolis-flux-localization}
Let \((W,\Pi)\) be as in \Cref{lem:R2-coriolis-bernoulli}, and let
\(\chi\in C_c^\infty(\mathbb R^3)\).  Then
\begin{equation}\label{eq:R2-local-coriolis-source}
   \int_{\mathbb R^3}\chi
   \left(
      -\Delta B+A\cdot\nabla B+|\omega|^2
   \right)\,dy
   =
   -2\int_{\mathbb R^3}(W\times\varpi)\cdot\nabla\chi\,dy .
\end{equation}
In particular the Coriolis contribution is supported only where the cutoff
varies.
\end{lemma}

\begin{proof}
By \eqref{eq:R2-Bernoulli-PDE},
\[
   -\Delta B+A\cdot\nabla B+|\omega|^2=2\varpi\cdot\omega .
\]
Because \(\varpi\) is constant and \(\omega=\nabla\times W\),
\[
   \varpi\cdot\omega
   =
   \nabla\cdot(W\times\varpi).
\]
Multiplying by \(\chi\), integrating over \(\mathbb R^3\), and integrating by
parts gives
\[
   2\int\chi\,\varpi\cdot\omega
   =
   2\int\chi\,\nabla\cdot(W\times\varpi)
   =
   -2\int(W\times\varpi)\cdot\nabla\chi ,
\]
which proves \eqref{eq:R2-local-coriolis-source}.
\end{proof}

\begin{lemma}[Non-small Coriolis flux creates a retained local concentration]
\label{lem:R2-annular-flux-active-concentration}
Let \(\chi\in C_c^\infty(\mathbb R^3)\), let
\(A_\chi:=\operatorname{supp}\nabla\chi\), and put
\[
   C_\chi:=\int_{A_\chi}|\nabla\chi|^{3/2}\,dy .
\]
Fix \(T\in(0,1]\), and let
\[
   \mathcal T_{\chi,T}
   :=
   \{(y,\tau): -T<\tau<0,\ Q(\tau)^Ty\in A_\chi\}
\]
be the swept annular tube in the original centered variables.  Cover
\(\mathcal T_{\chi,T}\) by finitely many unit terminal cylinders
\[
   \mathcal T_{\chi,T}\subset \bigcup_{m=1}^{N_{\chi,T}}Q_1(z_m).
\]
Let \(\eps_{\rm sd}>0\) be the local concentration threshold used in
\Cref{def:terminal-measures}.  If \(\varpi\ne0\) and
\begin{equation}\label{eq:R2-flux-active-threshold}
   \left|
      2\int_{\mathbb R^3}(W\times\varpi)\cdot\nabla\chi\,dy
   \right|
   \ge
   2|\varpi|\,C_\chi^{2/3}
      \left(\frac{N_{\chi,T}\eps_{\rm sd}}{T}\right)^{1/3},
\end{equation}
then at least one cylinder \(Q_1(z_m)\) is \(\eps_{\rm sd}\)-concentrating for the original
centered solution \(V\).
\end{lemma}

\begin{proof}
Write
\[
   \delta_\chi:=
   \left|
      2\int_{\mathbb R^3}(W\times\varpi)\cdot\nabla\chi\,dy
   \right| .
\]
Holder's inequality gives
\[
\begin{aligned}
   \delta_\chi
   &\le
   2|\varpi|\int_{A_\chi}|W|\,|\nabla\chi|\,dy  \\
   &\le
   2|\varpi|
   \left(\int_{A_\chi}|W|^3\,dy\right)^{1/3}
   C_\chi^{2/3}.
\end{aligned}
\]
Thus \eqref{eq:R2-flux-active-threshold} implies
\[
   \int_{A_\chi}|W|^3\,dy\ge \frac{N_{\chi,T}\eps_{\rm sd}}{T}.
\]
Since \(V(y,\tau)=Q(\tau)W(Q(\tau)^Ty)\) and \(Q(\tau)\) is an isometry,
\[
\begin{aligned}
   \iint_{\mathcal T_{\chi,T}} |V(y,\tau)|^3\,dy\,d\tau
   &=
   \int_{-T}^{0}\int_{Q(\tau)A_\chi}
      |W(Q(\tau)^Ty)|^3\,dy\,d\tau \\
   &=
   T\int_{A_\chi}|W(z)|^3\,dz
   \ge N_{\chi,T}\eps_{\rm sd}.
\end{aligned}
\]
Because the tube is covered by \(N_{\chi,T}\) unit terminal cylinders, one of
them satisfies
\[
   \iint_{Q_1(z_m)} |V|^3\,dy\,d\tau\ge \eps_{\rm sd}.
\]
The pressure contribution to the local CKN measure is nonnegative after the
allowed pressure gauge.  Hence
\[
   \mu(Q_1(z_m))\ge \eps_{\rm sd},
\]
so the corresponding recentering sequence is \(\eps_{\rm sd}\)-concentrating in the sense of
\Cref{def:terminal-measures}.
\end{proof}

\begin{lemma}[Local Coriolis-flux dichotomy]
\label{lem:R2-flux-routing-dichotomy}
Fix \(\chi\in C_c^\infty(\mathbb R^3)\) and \(T\in(0,1]\).  Along every
terminal extraction of a rotational relative equilibrium, exactly one of the
following local alternatives occurs:
\begin{enumerate}[label=(\roman*)]
   \item \(\varpi\ne0\), the flux across
   \(A_\chi=\operatorname{supp}\nabla\chi\) satisfies the threshold
   \eqref{eq:R2-flux-active-threshold}, and the swept annular tube contains a
   unit-scale concentration;
   \item either \(\varpi=0\), or the threshold
   \eqref{eq:R2-flux-active-threshold} fails; then the flux is a bounded
   compactly supported functional of \(W\) on \(A_\chi\), passes to terminal
   limits by local \(L^3\) convergence, and creates no terminal alternative other than
   the single-profile, finite separated-profile, exterior-escape,
   diffuse-exterior-concentration, local-vanishing, or
   local-compactness-failure alternatives of
   \Cref{thm:terminal-exhaustion-main}.
\end{enumerate}
\end{lemma}

\begin{proof}
Alternative \emph{(i)} is precisely
\Cref{lem:R2-annular-flux-active-concentration}.  If \(\varpi=0\), then the flux
vanishes identically.  If \(\varpi\ne0\) and
\eqref{eq:R2-flux-active-threshold} fails, the flux is bounded by the right hand
side of that threshold.  In both cases the functional
\[
   W\longmapsto -2\int_{\mathbb R^3}(W\times\varpi)\cdot\nabla\chi\,dy
\]
depends only on \(W\) on the compact annular set \(A_\chi\).  Under the local
Seregin compactness supplied by Paper I, velocities converge strongly in
\(L^3_{\rm loc}\); since \(\nabla\chi\in L^{3/2}\) and \(\varpi\) is fixed, the
functional is continuous along every terminal subsequence:
\[
   \int (W_n\times\varpi)\cdot\nabla\chi
   \longrightarrow
   \int (W_\infty\times\varpi)\cdot\nabla\chi .
\]
Thus a subthreshold flux is retained only as a local term in the limiting
localized Bernoulli identity \eqref{eq:R2-local-coriolis-source}.  It is not a
new location of local velocity mass.  Every actual location of compact velocity
\(L^3\)-mass is still one of the terminal local concentration alternatives listed in
\Cref{thm:terminal-exhaustion-main}.
\end{proof}

\begin{remark}[Why this is not a global Liouville argument]
\label{rem:R2-no-global-liouville}
For \(\Omega=0\), the Coriolis flux in
\Cref{lem:R2-coriolis-flux-localization} vanishes.  For \(\Omega\ne0\), the
same lemma does not give a sign-definite maximum principle for \(B\).  It gives
only the local dichotomy in \Cref{lem:R2-flux-routing-dichotomy}: either a
chosen annulus already contains a recentered local concentration, or the
annular flux is a subthreshold compact functional which passes to the local
terminal limit.  No tail decay, whole-space \(L^p\) bound, Gaussian energy
estimate, or global pressure normalization is used.
\end{remark}

\subsection{Local closure of the rotational relative-equilibrium class}

\begin{theorem}[Stratified exclusion of rotational relative equilibria]
\label{thm:R2-exclusion-safe}
The ordered rotational relative-equilibrium residual class is empty:
\begin{equation}\label{eq:R2-empty-safe}
   \Crot(\mathcal S;I,J)=\varnothing .
\end{equation}
\end{theorem}

\begin{proof}
Suppose, for contradiction, that
\(V\in\Crot(\mathcal S;I,J)\).  By
\Cref{def:R2-relative-profile-safe} and
\Cref{lem:co-rotating-equation-safe}, there are
\(\Omega\in\mathfrak{so}(3)\) and a bounded smooth divergence-free profile
\((W,\Pi)\) satisfying the Coriolis--Leray equation
\eqref{eq:R2-coriolis-profile-safe}.  The centered representative \(V\) is a
normalized Seregin ancient limit, so by
\Cref{hyp:base-seregin-package} it carries a fixed positive compact
velocity \(L^3\)-mass.

Choose a countable family of smooth cutoffs which is cofinal among compact
annuli in the terminal variables.  For each cutoff and each rational
\(T\in(0,1]\), apply the local dichotomy
\Cref{lem:R2-flux-routing-dichotomy}.  A superthreshold Coriolis flux produces
a unit-scale concentration and is therefore already part of the local
concentration decomposition.  A subthreshold Coriolis flux is a compactly
supported functional
of \(W\), continuous under the local convergence supplied by Paper I, and hence
remains inside the localized limiting equation; it carries no independent CKN
mass and creates no additional terminal alternative.

We now apply the terminal local concentration decomposition to the inherited compact
density.  Loss of local compactness, local vanishing, exterior escape, and
diffuse exterior concentration cannot be the sole locations of this density by
\Cref{thm:terminal-nonactive-discharge}.  By the preceding paragraph, the
Coriolis flux has already led either to a local concentration or to a
compact local term in the selected terminal equation.  Thus, after local
concentration extraction, the retained density is located either in a finite separated family
of retained concentration profiles or in a single retained concentration
terminal profile.

The finite separated-family alternative is impossible by
\Cref{thm:no-finite-separated-profile-family}.  In the single-profile case,
\Cref{prop:no-multi-implies-atomic} makes the terminal profile terminally
indecomposable, and \Cref{lem:no-atomic-active} excludes retained
indecomposable concentration profiles.  Both possible locations of the inherited
compact density are impossible.  This
contradicts the normalization of \(V\), and therefore
\(\Crot(\mathcal S;I,J)=\varnothing\).
\end{proof}

\begin{proposition}[Already-closed integrable subcase]
\label{prop:R2-integrable-subcase}
Let \(V\) be a rotational relative equilibrium with co-rotating profile \(W\).
If
\begin{equation}\label{eq:R2-L3-subcase}
   W\in L^3(\mathbb R^3),
\end{equation}
then \(V\equiv0\).  More generally, the same conclusion holds under any
condition which places the physical representative in
\(L^\infty_tL^3_x\) on its ancient time interval.
\end{proposition}

\begin{proof}
Rotations and the Navier--Stokes critical scaling preserve the \(L^3\) norm.
Thus \(W\in L^3\) implies
\[
   \sup_{\tau\in\mathbb R}\|V(\cdot,\tau)\|_{L^3(\mathbb R^3)}<\infty .
\]
Equivalently, the physical representative lies in the critical class
\(L^\infty_tL^3_x\).  The endpoint \(L^3\) class is one of the already closed
alternatives in Paper I, by the Escauriaza--Seregin--\v Sver\'ak endpoint
regularity theorem \cite{ESS2003} and the stationary
Ne\v cas--R\r u\v zi\v cka--\v Sver\'ak Liouville theorem \cite{NRS1996}.
Therefore
\(V\equiv0\).  A zero profile has zero local velocity mass, so no such profile
can belong to the residual
nonvanishing class.
\end{proof}


\section{The degenerate stationary-hull residual class}
\label{sec:R3-degenerate-stationary-hull}

We now isolate the third refined residual class.  This class is designed to
capture Type~I ancient limits whose time-translation hull contains a stationary
profile with retained local concentration lying on a non-isolated family of
stationary centered profiles.  The role of the stationary-family condition is to separate genuinely degenerate
stationary profile families from isolated stationary profiles and from the
symmetry-generated classes already removed in \(\Cax\) and \(\Crot\).

A small point of formulation is important.  It is not enough to assume that the
hull of a residual ancient solution contains some stationary-family element.  A
stationary element may occur as a limit without retained local concentration and
therefore need not retain the local CKN concentration inherited from the
original singular point.  Such a hull element should not be placed in a class
that is to be closed by a
stationary-profile exclusion.  Therefore the class below is defined using
\emph{retained-concentration occurrence}: the stationary-family profile must
itself occur as a Seregin limit carrying a positive local CKN concentration.

Throughout this section the centered equation is
\begin{equation}\label{eq:centered-R3}
   \partial_\tau V-\Delta V+\frac12 y\cdot\nabla V+\frac12 V
   +(V\cdot\nabla)V+\nabla P=0,
   \qquad \nabla\cdot V=0 .
\end{equation}
Its projected form is
\begin{equation}\label{eq:projected-centered-R3}
   \partial_\tau V=\mathcal F(V),
   \qquad
   \mathcal F(W)
   :=
   \mathbb P\left[
      \Delta W-\frac12(W+y\cdot\nabla W)-(W\cdot\nabla)W
   \right],
\end{equation}
where \(\mathbb P\) is the whole-space Leray projector.

\subsection{Stationary families and retained-concentration hull limits}

\begin{definition}[Admissible stationary phase space]
\label{def:R3-stationary-phase-space}
An admissible stationary phase space is a Banach space \(\mathfrak B\) of
solenoidal vector fields on \(\mathbb R^3\) such that:
\begin{enumerate}[label=(\roman*)]
   \item \(\mathfrak B\hookrightarrow C^{2,\alpha}_{\loc}(\mathbb R^3)
      \cap L^\infty(\mathbb R^3)\) continuously for some
      \(\alpha\in(0,1)\);
   \item the map \(\mathcal F\) in \eqref{eq:projected-centered-R3} is
      well-defined on an open subset \(\mathcal U\subset\mathfrak B\) and maps
      \(\mathcal U\) into a Banach space \(\mathfrak Y\);
   \item \(\mathcal F:\mathcal U\to\mathfrak Y\) is \(C^1\).
\end{enumerate}
\end{definition}

\begin{definition}[Nontrivial stationary family]
\label{def:R3-stationary-family}
Let \(W\in\mathcal U\subset\mathfrak B\) satisfy \(\mathcal F(W)=0\).  We say
that \(W\) lies on an admissible nontrivial stationary family if there exist
\(\varepsilon>0\) and a \(C^1\) curve
\[
   a\mapsto W_a\in\mathcal U,
   \qquad |a|<\varepsilon,
\]
such that
\begin{equation}\label{eq:R3-stationary-family}
   W_0=W,
   \qquad
   \mathcal F(W_a)=0 \quad (|a|<\varepsilon),
\end{equation}
and such that the family is not locally generated only by the action of the
rigid rotation group.  Equivalently, after possibly reparametrizing the curve,
one may assume that
\begin{equation}\label{eq:R3-nontrivial-tangent}
   \dot W_0\notin
   \{\Omega W-(\Omega y)\cdot\nabla W:\Omega\in\mathfrak{so}(3)\}
\end{equation}
whenever the tangent \(\dot W_0\) exists and is nonzero.
\end{definition}

\begin{definition}[Retained local CKN concentration]
\label{def:R3-active-profile}
Let \((W,\Pi)\) be a stationary centered profile.  We say that \((W,\Pi)\) is
\(\varepsilon_*\)-concentrating if there exist a compact parabolic cylinder
\(Q_r(z_0)\) and an admissible pressure gauge \(c_{r,z_0}\) such that
\begin{equation}\label{eq:R3-active-ckn}
   \iint_{Q_r(z_0)}
   \left(|W|^3+|\Pi-c_{r,z_0}|^{3/2}\right)\,dy\,d\tau
   \ge \varepsilon_* .
\end{equation}
For a stationary profile the integrand is independent of \(\tau\), but the
parabolic-cylinder notation is retained to match the CKN normalization used for
Seregin limits.
\end{definition}

\begin{definition}[Degenerate stationary-hull residual class with retained concentration]
\label{def:R3-active-class}
Fix an admissible stationary phase space \(\mathfrak B\).  The degenerate
stationary-hull residual class with retained local concentration
\[
   \Cstat(\mathcal S;I,J;\mathfrak B)
\]
consists of those \(V\in\mathcal R(\mathcal S;I,J)\), after the ordered
removal of \(\Cax\) and \(\Crot\), for which there exists a
sequence \(s_k\in\mathbb R\) and a stationary profile \((W,\Pi)\) such that
\begin{equation}\label{eq:R3-hull-convergence}
   (V(\cdot,\cdot+s_k),P(\cdot,\cdot+s_k)-c_k(\cdot))
   \longrightarrow (W,\Pi)
\end{equation}
locally smoothly for the velocity and locally in \(L^{3/2}\) for the pressure
modulo admissible gauges, and such that:
\begin{enumerate}[label=(\roman*)]
   \item \(W\in\mathfrak B\) and \(\mathcal F(W)=0\);
   \item \(W\) lies on an admissible nontrivial stationary family in the sense
      of \Cref{def:R3-stationary-family};
   \item \((W,\Pi)\) is \(\varepsilon_*\)-concentrating in the sense of
      \Cref{def:R3-active-profile};
   \item \((W,\Pi)\) is not in any previously closed stationary, structured,
      axisymmetric bounded-circulation, or rotational relative-equilibrium
      class.
\end{enumerate}
In the ordered residual decomposition one keeps this class of stationary-hull
limits with retained concentration as the degenerate stationary-hull
alternative.  The hull cases
without retained local concentration
should be assigned to the generic terminal residual class, where they are
handled by the terminal local concentration decomposition rather than by a
stationary-profile exclusion.
\end{definition}

\begin{remark}[Why retained concentration is necessary]
\label{rem:R3-active-necessary}
The condition \eqref{eq:R3-active-ckn} prevents the following logical error.  A
nonstationary ancient solution may have a stationary profile in its hull with
zero local CKN mass on the relevant compact cylinders.  Excluding that stationary
profile would not exclude the original ancient solution.  The
stationary-profile exclusion used below applies only to stationary profiles
that themselves occur as nontrivial Seregin limits with positive local
concentration.
\end{remark}

\subsection{Hull limits are again Seregin limits}

The next lemma is the basic compatibility between time-translation hulls in
centered variables and rescaling sequences in physical variables.  It is the
reason that a stationary profile in the hull which retains local concentration
may be treated as a blow-up profile of the original solution, not merely as a
limit inside the already extracted ancient solution.

\begin{lemma}[Scale reselection for centered hull limits]
\label{lem:R3-hull-is-seregin}
Let \((V,P)\) be a centered Type~I Seregin limit obtained from a finite-energy
suitable weak solution \((u,p)\) at a singular point \(z_*=(x_*,t_*)\).  Suppose
that, for some sequence \(s_k\in\mathbb R\),
\[
   (V(\cdot,\cdot+s_k),P(\cdot,\cdot+s_k)-c_k(\cdot))
   \to (W,\Pi)
\]
locally in the Seregin topology.  Then, after passing to a diagonal subsequence
and reselecting the physical blow-up scales, \((W,\Pi)\) is itself a centered
Type~I Seregin limit of \((u,p)\) at \(z_*\).  If the convergence to \((W,\Pi)\)
retains a positive local CKN lower bound, then \((W,\Pi)\) is a nontrivial
Seregin limit with the same type of inherited nonvanishing condition.
\end{lemma}

\begin{proof}
Let \(\lambda_n\downarrow0\) be a sequence of scales producing \((V,P)\).  In
centered variables, a time translation by \(s\) is exactly a change of the
physical blow-up scale by a factor depending only on \(s\).  With the convention
used in \eqref{eq:centered-R3}, write this adjusted scale as
\begin{equation}\label{eq:R3-adjusted-scale}
   \lambda_{n,s}:=e^{-s/2}\lambda_n .
\end{equation}
The rescaled fields built from \((u,p)\) at scale \(\lambda_{n,s}\) are the same
as the fields at scale \(\lambda_n\), translated by \(s\) in the centered time
variable, up to the corresponding pressure gauge.

Fix an exhaustion of compact cylinders \(K_m\subset\mathbb R^3\times\mathbb R\).
For each \(k\), choose \(n(k)\) so large that \(\lambda_{n(k),s_k}\to0\) and the
rescaled fields at scale \(\lambda_{n(k)}\), shifted by \(s_k\), are within
\(1/k\) of \((V(\cdot,\cdot+s_k),P(\cdot,\cdot+s_k)-c_k)\) on \(K_k\) in the
local Seregin topology.  This is possible because \(\lambda_n\downarrow0\) and
because the original sequence converges locally to \((V,P)\).

By the assumed hull convergence, the shifted limit
\((V(\cdot,\cdot+s_k),P(\cdot,\cdot+s_k)-c_k)\) converges to \((W,\Pi)\) on each
fixed compact cylinder.  The diagonal sequence of physical rescalings with
scales \(\lambda_{n(k),s_k}\) therefore converges locally to \((W,\Pi)\).  Hence
\((W,\Pi)\) is again a centered Type~I Seregin limit at \(z_*\).  The assertion
about nonvanishing follows from lower semicontinuity and local convergence of
the CKN quantities after the prescribed pressure gauges.
\end{proof}

\subsection{Elementary consequences of stationary-family degeneracy}

The stationary-family condition has two immediate PDE consequences.  First, a differentiable
nontrivial family produces a nontrivial element of the kernel of the linearized
stationary Leray operator.  Second, if an ancient trajectory is exactly
constrained to a stationary family, then it is stationary.

\begin{lemma}[Stationary family implies linearized degeneracy]
\label{lem:R3-family-kernel}
Let \(a\mapsto W_a\in\mathfrak B\) be a \(C^1\) stationary family satisfying
\eqref{eq:R3-stationary-family}.  Then
\begin{equation}\label{eq:R3-kernel}
   D\mathcal F(W_0)\dot W_0=0 .
\end{equation}
In particular, if \(\dot W_0\ne0\) modulo infinitesimal rotations, the
linearized stationary Leray operator has a genuine non-symmetry kernel element.
\end{lemma}

\begin{proof}
Differentiate \(\mathcal F(W_a)=0\) at \(a=0\).  Since \(\mathcal F\) is \(C^1\)
on \(\mathfrak B\), this gives \eqref{eq:R3-kernel}.  The final statement is
exactly the nontriviality condition in \Cref{def:R3-stationary-family}.
\end{proof}

\begin{lemma}[Trajectories constrained to a stationary family are stationary]
\label{lem:R3-family-constrained-stationary}
Let \(a\mapsto W_a\in\mathfrak B\) be a \(C^1\) stationary family with
\(\dot W_a\ne0\) on an interval.  Suppose that a smooth solution of
\eqref{eq:projected-centered-R3} has the form
\begin{equation}\label{eq:R3-family-trajectory}
   V(\tau)=W_{a(\tau)}
\end{equation}
for a \(C^1\) function \(a(\tau)\).  Then \(a\) is constant, and hence
\(V\) is stationary.
\end{lemma}

\begin{proof}
Since \(\mathcal F(W_a)=0\) for every \(a\), equation
\eqref{eq:projected-centered-R3} gives
\[
   a'(\tau)\dot W_{a(\tau)}=\partial_\tau V(\tau)=\mathcal F(W_{a(\tau)})=0 .
\]
Because \(\dot W_{a(\tau)}\ne0\), we obtain \(a'(\tau)=0\).  Thus \(a\) is
constant on the interval.
\end{proof}

\begin{corollary}[Exact stationary-family subcase]
\label{cor:R3-exact-family-subcase}
If an element of the residual class is itself an exact trajectory on an
admissible stationary family, then it is a stationary centered profile.  Hence
it is excluded whenever it satisfies the stationary \(L^3\) Liouville
criterion, or more generally whenever the stationary profile belongs to one of
the previously closed stationary or structured classes.
\end{corollary}

\begin{proof}
Let \(V(\tau)=W_{a(\tau)}\) be such an exact stationary-family trajectory.  By
\Cref{lem:R3-family-constrained-stationary}, \(a'(\tau)=0\) on every interval on
which the stationary-family parametrization is nondegenerate.  Hence \(a\) is constant and
\(V\) is a single stationary centered profile.  If that stationary profile
satisfies the stationary \(L^3\) Liouville criterion or lies in one of the
previously closed structured or ordered lower residual classes, the corresponding
prior exclusion removes it from the residual class.
\end{proof}

\subsection{Stationary-hull profiles which occur as Seregin limits}

The remaining issue is not the internal dynamics along a stationary family; exact
stationary-family
motion is stationary by \Cref{lem:R3-family-constrained-stationary}.  The real
issue is occurrence as a blow-up limit: can a degenerate stationary-family
profile occur as a Seregin limit of a finite-energy suitable weak solution
while retaining local concentration?  The following statement is the precise
Seregin limit occurrence statement
required to close the degenerate stationary-hull class with retained local
concentration.  It is proved in the following section by the terminal local
concentration theorem.

\begin{theorem}[No degenerate stationary-family concentration profile occurs as a Seregin limit]
\label{thm:R3-no-attainable-degenerate-family}
Fix an admissible stationary phase space \(\mathfrak B\).  Let \((W,\Pi)\) be a
stationary centered profile satisfying:
\begin{enumerate}[label=(\roman*)]
   \item \(W\in\mathfrak B\), \(\mathcal F(W)=0\), and \(W\) lies on an
      admissible nontrivial stationary family;
   \item \((W,\Pi)\) is a centered Type~I Seregin limit of a finite-energy
      suitable weak solution at a singular point;
   \item \((W,\Pi)\) is \(\varepsilon_*\)-concentrating;
   \item \((W,\Pi)\) is not in the small, stationary \(L^3\), tight, closed
      structured/decay, axisymmetric bounded-circulation, or rotational
      relative-equilibrium classes already removed earlier in the stratification.
\end{enumerate}
Then no such \((W,\Pi)\) exists.
\end{theorem}

\begin{proof}
This is exactly \Cref{cor:R3S-no-active-degenerate-family}, proved below from
the local terminal stratification theorem.
\end{proof}

\begin{remark}[Interpretation of the Seregin limit occurrence theorem]
\label{rem:R3-input-interpretation}
\Cref{thm:R3-no-attainable-degenerate-family} excludes occurrence
as a Type~I Seregin limit.  It does not assert that degenerate stationary profiles cannot
exist as elliptic solutions of the stationary centered equation.  It asserts
only that such profiles cannot occur as Type~I blow-up limits with retained
local concentration after all previously closed classes have been removed.  This
is the correct formulation for the residual problem: the elliptic family may
exist, but it cannot be reached by the local blow-up procedure coming from a
finite-energy suitable weak solution.
\end{remark}

\subsection{Closure of the degenerate stationary-hull concentration class}

\begin{theorem}[Exclusion of the degenerate stationary-hull concentration class]
\label{thm:R3-exclusion}
\begin{equation}\label{eq:R3-empty}
   \Cstat(\mathcal S;I,J;\mathfrak B)=\varnothing .
\end{equation}
Consequently the ordered third residual class is empty, provided it is defined
with the retained-concentration Seregin limit condition in
\Cref{def:R3-active-class}.
\end{theorem}

\begin{proof}
Suppose, for contradiction, that
\(V\in\Cstat(\mathcal S;I,J;\mathfrak B)\).  By
\Cref{def:R3-active-class}, there is a stationary profile \((W,\Pi)\) in the
centered time-translation hull of \((V,P)\), and \((W,\Pi)\) is
\(\varepsilon_*\)-concentrating.  The profile \(W\) belongs to \(\mathfrak B\), satisfies
\(\mathcal F(W)=0\), and lies on an admissible nontrivial stationary family.
Moreover, by the ordered definition of \(\Cstat\), it is not
in any of the previously closed lower classes.

By \Cref{lem:R3-hull-is-seregin}, the hull profile \((W,\Pi)\) is itself a
centered Type~I Seregin limit of the original finite-energy suitable weak
solution at the same singular point, after reselecting scales.  Since local
concentration is retained, \((W,\Pi)\) satisfies every condition in
\Cref{thm:R3-no-attainable-degenerate-family}.  This contradicts the
Seregin limit occurrence theorem.
Therefore no such \(V\) exists, proving \eqref{eq:R3-empty}.
\end{proof}

\begin{proposition}[Already-closed integrable stationary-family subcase]
\label{prop:R3-L3-subcase}
Let \((W,\Pi)\) be a stationary centered profile lying on an admissible
stationary family.  If
\begin{equation}\label{eq:R3-L3}
   W\in L^3(\mathbb R^3),
\end{equation}
then \(W\equiv0\).  Hence no retained-concentration hull profile satisfying
\eqref{eq:R3-L3} can occur in \(\Cstat\).
\end{proposition}

\begin{proof}
The stationary physical representative associated with \(W\) is a backward
self-similar Navier--Stokes solution.  The assumption \(W\in L^3(\mathbb R^3)\)
places it in the stationary \(L^3\) class already excluded by the classical
Ne\v cas--R\r u\v zi\v cka--\v Sver\'ak Liouville theorem \cite{NRS1996},
equivalently by the endpoint ancient Liouville theorem recorded in Paper I and
cited below as \Cref{thm:AB-main}.  Hence \(W\equiv0\).  A
zero profile has zero pressure-normalized local CKN mass after choosing the
admissible pressure gauge, so it cannot be \(\varepsilon_*\)-concentrating.
\end{proof}

\begin{remark}[Closure after terminal local concentration analysis]
\label{rem:R3-final-obligation}
The preceding argument reduces the stationary-hull class with retained local
concentration from a time-translation hull question to the question whether a
stationary centered profile in a nontrivial stationary family can occur as a
Seregin limit.
\Cref{thm:R3-no-attainable-degenerate-family} is supplied by the terminal local
concentration analysis in the next section.
\end{remark}

\begin{remark}[Recommended statement for the residual decomposition]
\label{rem:R3-decomposition-replacement}
For the residual decomposition, the third exclusion line should be stated as follows:
\[
\begin{gathered}
   \text{stationary-hull residual with retained concentration}\\
   \xrightarrow{\;\text{stationary-family Seregin limit exclusion}\;}
   \varnothing .
\end{gathered}
\]
The precise class is the
\emph{degenerate stationary-hull residual class with retained local
concentration}.
\end{remark}


% -----------------------------------------------------------------------------
\section{Exclusion of degenerate stationary-hull profiles with retained concentration}
\label{sec:R3-terminal-stratification-closure}

We close the remaining analytic statement in the degenerate stationary-hull
class by the terminal local concentration theorem proved in the main
paper.  The argument is not specific to the existence of a
non-isolated stationary family.  Once a stationary-family element occurs as
a Type~I Seregin limit carrying retained local CKN concentration, it is a
stationary terminal concentration profile.  Such profiles are excluded by the
same local terminal concentration analysis and separated-profile exclusion used
for the final residual class.  No global coercive
estimate for the original solution is used.

Throughout this section, a centered Type~I profile satisfies
\begin{equation}\label{eq:R3S-centered}
   \partial_\tau U-\Delta U+\frac12 y\cdot\nabla U+\frac12 U
   +(U\cdot\nabla)U+\nabla \Pi=0,
   \qquad \nabla\cdot U=0
\end{equation}
on \(\mathbb R^3\times\mathbb R\), with the usual local suitability and
pressure-gauge conventions.

\subsection{Terminal local concentration theorem}

We use the following specialization of the terminal local concentration theorem
proved below.  It is stated here exactly in the form needed for the
stationary-hull closure with retained local concentration.

\begin{theorem}[Terminal local concentration exclusion theorem]
\label{thm:R3S-terminal-package}
Let \((U,\Pi)\) be a bounded centered Type~I Seregin profile satisfying local
suitability and carrying a retained local CKN concentration, i.e. for some
compact parabolic cylinder \(Q\) and some \(\varepsilon_*>0\),
\begin{equation}\label{eq:R3S-retained-activity}
   \iint_Q\bigl(|U|^3+|\Pi-c_Q(\tau)|^{3/2}\bigr)\,dy\,d\tau
   \ge \varepsilon_* .
\end{equation}
Assume that \((U,\Pi)\) is outside the previously closed small-amplitude,
stationary \(L^3\), tight, closed structured/decay, axisymmetric
bounded-circulation, and rotational relative-equilibrium classes.  Then the
following hold.
\begin{enumerate}[label=(\roman*)]
   \item Every terminal extraction of \((U,\Pi)\) belongs to one of the local
   alternatives in \Cref{thm:terminal-exhaustion-main}: a single retained
   concentration profile, local vanishing after removal of concentrating
   sequences, exterior escape of the concentrating sequence, diffuse exterior
   concentration, loss of local compactness, or a finite separated family of
   retained concentration profiles.
   \item Local vanishing, exterior escape, diffuse exterior concentration, and
   loss of local compactness cannot be the sole locations of the retained
   compact CKN concentration
   \eqref{eq:R3S-retained-activity}.
   \item Finite separated families of retained concentration profiles are
   impossible.
   \item Any remaining single retained concentration profile is terminally
   indecomposable: it has no retained concentrating tail limit, no diffuse
   exterior concentration, and no loss of local compactness along escaping
   recenterings.
   \item Every retained terminally indecomposable concentration profile has a
   backward sequence of finite critical norm: there exist
   \(\tau_k\to-\infty\) such that
   \begin{equation}\label{eq:R3S-atomic-L3-input}
      \sup_k \|U(\cdot,\tau_k)\|_{L^3(\mathbb R^3)}<\infty .
   \end{equation}
   \item Retained terminal concentration profiles inherit mildness under the
   physical pullback.  Therefore the Albritton--Barker endpoint Liouville theorem
   applies to any profile satisfying \eqref{eq:R3S-atomic-L3-input}.
\end{enumerate}
\end{theorem}

\begin{proof}
This is \Cref{thm:terminal-local-stratification-package} applied to
\((U,\Pi)\).  The local alternatives in (i) are those of
\Cref{thm:terminal-exhaustion-main}.  Item (ii) is the exclusion of local
vanishing, exterior escape, diffuse exterior concentration, and local
compactness failure in \Cref{thm:terminal-nonactive-discharge}.  Item (iii) is
\Cref{thm:no-finite-separated-profile-family}.  Item (iv) is
\Cref{prop:no-multi-implies-atomic}.  Item (v) is
\Cref{lem:atomic-sequence-L3}, and item (vi) is
\Cref{hyp:mildness-inheritance-main,thm:AB-main}.  These results are proved in
the present paper from Paper 0, Paper I, and local concentration analysis;
no global bound is assumed.
\end{proof}

\begin{remark}[Local nature of the theorem]
\label{rem:R3S-local-nature}
The theorem above is a local concentration statement.  It uses the local
suitable-weak-solution compactness and pressure normalization developed from
the Caffarelli--Kohn--Nirenberg theory \cite{CKN1982}, together with the
Paper 0 local concentration input and the Paper I Seregin compactness theorem.
It then uses terminal profile decomposition, exclusion of local vanishing,
diffuse exterior concentration, and local compactness failure, and the finite
separated-profile theorem.  It does not
assert or require a global a priori bound for the original physical solution.
\end{remark}

For completeness, and to make the dependence on local estimates explicit, we
record the elementary part of the argument which converts terminal
indecomposability into sequence-\(L^3\) control.

\begin{lemma}[Indecomposable local profiles have backward sequence-\(L^3\) control]
\label{lem:R3S-atomic-L3}
Let \((U,\Pi)\) be a bounded smooth ancient solution of
\eqref{eq:R3S-centered}.  Suppose that \((U,\Pi)\) has retained local
concentration and is terminally indecomposable: it has no retained
concentrating tail limit, no
diffuse exterior concentration, and no escaping recentering sequence along
which local compactness, suitability, or pressure control fails.  Then there
exists a sequence \(\tau_k\to-\infty\) such that
\begin{equation}\label{eq:R3S-atomic-L3}
   \sup_k \|U(\cdot,\tau_k)\|_{L^3(\mathbb R^3)}<\infty .
\end{equation}
\end{lemma}

\begin{proof}
We argue by contradiction.  Assume that no sequence satisfying
\eqref{eq:R3S-atomic-L3} exists.  Then
\begin{equation}\label{eq:R3S-L3-divergence}
   \liminf_{\tau\to-\infty}\|U(\cdot,\tau)\|_{L^3(\mathbb R^3)}=\infty .
\end{equation}
Choose \(\tau_n\to-\infty\) such that
\(
   \int_{\mathbb R^3}|U(y,\tau_n)|^3\,dy\to\infty .
\)
Since \(U\) is bounded, we may choose radii \(R_n\to\infty\) so that
\begin{equation}\label{eq:R3S-tail-mass}
   \int_{|y|>R_n}|U(y,\tau_n)|^3\,dy\ge 1 .
\end{equation}
Let
\begin{equation}\label{eq:R3S-unit-tail-sup}
   m_n:=\sup_{|x|>R_n}\int_{B_1(x)}|U(y,\tau_n)|^3\,dy .
\end{equation}
If \(\limsup_n m_n\ge\varepsilon_*\), then for a subsequence there exist
\(|x_n|>R_n\) such that
\begin{equation}\label{eq:R3S-active-tail-cube}
   \int_{B_1(x_n)}|U(y,\tau_n)|^3\,dy\ge \varepsilon_* .
\end{equation}
After recentering at \((x_n,\tau_n)\), either the local compactness and pressure
control required for a terminal profile fail, or a subsequence converges in the
terminal topology to a retained concentrating tail limit.  Both alternatives
contradict terminal indecomposability.

It remains to consider the case \(m_n<\varepsilon_*\) for all sufficiently large
\(n\).  Then \eqref{eq:R3S-tail-mass} says that critical mass persists outside
expanding balls, while \eqref{eq:R3S-unit-tail-sup} says that every unit-scale
exterior recentering remains below the retained-concentration threshold.  This
is exactly diffuse exterior concentration, again contradicting terminal
indecomposability.
Thus \eqref{eq:R3S-L3-divergence} is impossible, and
\eqref{eq:R3S-atomic-L3} follows.
\end{proof}

\subsection{Stationary concentration profiles are impossible}

\begin{theorem}[No stationary residual concentration profile occurs as a Seregin limit]
\label{thm:R3S-no-active-stationary-carrier}
Let \((W,\Pi)\) be a stationary solution of the centered equation,
\begin{equation}\label{eq:R3S-stationary-centered}
   -\Delta W+\frac12 y\cdot\nabla W+\frac12 W
   +(W\cdot\nabla)W+\nabla\Pi=0,
   \qquad \nabla\cdot W=0 .
\end{equation}
Assume that:
\begin{enumerate}[label=(\roman*)]
   \item \((W,\Pi)\) is obtained as a centered Type~I Seregin limit of a
   finite-energy suitable weak solution at a singular point;
   \item \((W,\Pi)\) has retained local concentration in the sense of
   \eqref{eq:R3S-retained-activity};
   \item \((W,\Pi)\) is outside all previously closed lower classes listed in
   \Cref{thm:R3S-terminal-package};
\end{enumerate}
Then no such \((W,\Pi)\) exists.
\end{theorem}

\begin{proof}
Suppose that such a stationary profile \((W,\Pi)\) exists.  Since it is a
centered Type~I Seregin limit and has retained local concentration, it is an admissible input
for the terminal local concentration theorem.  Apply \Cref{thm:R3S-terminal-package} to
\((W,\Pi)\).

Finite separated terminal concentration profiles are impossible by
\Cref{thm:R3S-terminal-package}.  Local vanishing, exterior escape, diffuse
exterior concentration, and local compactness failure cannot be the only
locations of the retained compact CKN concentration.  Hence the retained
concentration is represented by a single terminal concentration profile.  By
the indecomposability conclusion in \Cref{thm:R3S-terminal-package}, this
single profile is terminally indecomposable.

Because \((W,\Pi)\) is stationary in \(\tau\), the single terminal
concentration profile may be represented by \((W,\Pi)\) itself after the
terminal normalization.  Therefore
\Cref{lem:R3S-atomic-L3} gives a sequence \(\tau_k\to-\infty\) such that
\begin{equation}\label{eq:R3S-W-L3-sequence}
   \sup_k\|W\|_{L^3(\mathbb R^3)}
   =
   \sup_k\|W(\cdot,\tau_k)\|_{L^3(\mathbb R^3)}
   <\infty .
\end{equation}
Thus \(W\in L^3(\mathbb R^3)\).

Let
\begin{equation}\label{eq:R3S-physical-pullback}
   w(x,t)=(-t)^{-1/2}W\!\left(\frac{x}{\sqrt{-t}}\right),
   \qquad
   q(x,t)=(-t)^{-1}\Pi\!\left(\frac{x}{\sqrt{-t}}\right),
   \qquad t<0 .
\end{equation}
By \Cref{hyp:mildness-inheritance-main}, \((w,q)\) is a mild ancient
Navier--Stokes solution after any finite time shift.  The critical scaling gives
\begin{equation}\label{eq:R3S-critical-L3-scaling}
   \|w(\cdot,t)\|_{L^3(\mathbb R^3)}=\|W\|_{L^3(\mathbb R^3)}<\infty
   \qquad (t<0).
\end{equation}
Applying the Albritton--Barker endpoint Liouville theorem along any sequence
\(t_k\downarrow-\infty\), \Cref{thm:AB-main}, we obtain \(w\equiv0\), and hence
\(W\equiv0\).
This contradicts retained local concentration, since the pressure-normalized
local CKN quantity of the zero profile is zero.  Therefore no stationary
residual concentration profile can occur as such a Seregin limit.
\end{proof}

\subsection{Closure of the degenerate stationary-hull residual class}

\begin{corollary}[No degenerate stationary-family concentration profile occurs as a Seregin limit]
\label{cor:R3S-no-active-degenerate-family}
Fix an admissible stationary phase space \(\mathfrak B\).  There is no stationary
centered profile \((W,\Pi)\) satisfying all of the following conditions:
\begin{enumerate}[label=(\roman*)]
   \item \(W\in\mathfrak B\), \(W\) solves the stationary centered equation
   \eqref{eq:R3S-stationary-centered}, and \(W\) lies on an admissible
   nontrivial stationary family;
   \item \((W,\Pi)\) is a centered Type~I Seregin limit of a finite-energy
   suitable weak solution at a singular point;
   \item \((W,\Pi)\) has retained local concentration;
   \item \((W,\Pi)\) is outside the previously closed lower classes;
\end{enumerate}
\end{corollary}

\begin{proof}
The stationary-family condition in (i) is additional structure, but the exclusion in
\Cref{thm:R3S-no-active-stationary-carrier} applies to every stationary
residual profile carrying retained local concentration.  Hence it applies in
particular to stationary profiles lying in admissible nontrivial stationary
families.
\end{proof}

\begin{theorem}[Closure of the degenerate stationary-hull concentration class]
\label{thm:R3S-R3-empty}
Let \(\Cstat(\mathcal S;I,J;\mathfrak B)\) be the degenerate
stationary-hull residual class with retained local concentration.  Then
\begin{equation}\label{eq:R3S-R3-empty}
   \Cstat(\mathcal S;I,J;\mathfrak B)=\varnothing .
\end{equation}
\end{theorem}

\begin{proof}
Suppose, for contradiction, that
\(V\in\Cstat(\mathcal S;I,J;\mathfrak B)\).  By the
definition of the degenerate stationary-hull residual class with retained local
concentration, the centered translation hull of \((V,P)\) contains a stationary
profile \((W,\Pi)\) with retained local concentration
which lies in an admissible nontrivial stationary family and is outside the
previously closed lower classes.  The scale-reselection lemma for centered hull
limits shows that \((W,\Pi)\) is itself a centered Type~I Seregin limit of the
original finite-energy suitable weak solution at the same singular point.

Thus \((W,\Pi)\) satisfies the hypotheses of
\Cref{cor:R3S-no-active-degenerate-family}.  This is impossible.  Hence no such
\(V\) exists, and \eqref{eq:R3S-R3-empty} follows.
\end{proof}

\begin{remark}[How this enters the residual decomposition]
\label{rem:R3S-decomposition}
The third residual line should be recorded as a theorem closed by the terminal
local concentration analysis:
\[
\begin{gathered}
   \text{degenerate stationary-hull residual with retained concentration}\\
   \Longrightarrow
   \text{stationary profile with retained concentration}
   \Longrightarrow
   \varnothing .
\end{gathered}
\]
The second implication is proved by terminal local concentration analysis:
decomposable profiles produce either an excluded separated concentration
family, diffuse exterior concentration, or loss of local compactness;
indecomposable profiles have backward sequence-\(L^3\) control and are ruled
out by the endpoint ancient Liouville theorem.
\end{remark}


\section{Tail geometry of the generic terminal residual}
\label{sec:R4-tail-geometry}

We now pass to the generic terminal residual class.  The preliminary point is that
the centered equation is not invariant under raw fixed spatial translations.
The correct far-field recenterings are self-similar covariant translations.

\begin{definition}[Covariant translations]
\label{def:covariant-translations}
For \(a\in\R^3\), define
\[
   (T_aV)(y,\tau):=V(y+e^{\tau/2}a,\tau),
   \qquad
   (T_aP)(y,\tau):=P(y+e^{\tau/2}a,\tau).
\]
\end{definition}

\begin{lemma}[Covariant translation invariance]
\label{lem:covariant-translation-invariance}
If \((V,P)\) solves \eqref{eq:centered-main}, then
\((T_aV,T_aP)\) also solves \eqref{eq:centered-main}.
\end{lemma}

\begin{proof}
Set \(z=y+e^{\tau/2}a\).  Then
\[
   \partial_\tau(T_aV)(y,\tau)
   =
   (\partial_\tau V)(z,\tau)
   +\frac12e^{\tau/2}a\cdot\nabla V(z,\tau).
\]
Moreover,
\[
   \frac12 y\cdot\nabla(T_aV)(y,\tau)
   =
   \frac12 z\cdot\nabla V(z,\tau)
   -\frac12e^{\tau/2}a\cdot\nabla V(z,\tau).
\]
The two additional transport terms cancel in
\(\partial_\tau T_aV+\frac12y\cdot\nabla T_aV\).  The Laplacian, the
zeroth-order term, the convection term, the pressure gradient, and the
divergence constraint are translation covariant.  Substitution into
\eqref{eq:centered-main} proves the claim.
\end{proof}

\begin{remark}[Raw translates create a spurious drift]
\label{rem:raw-translates}
If one sets \(V_n(y,\tau)=V(y+x_n,\tau)\) with fixed \(x_n\), then the centered
equation acquires the additional transport term
\[
   -\frac12 x_n\cdot\nabla V_n .
\]
Thus fixed translates do not define a compact invariant far-field hull for the
renormalized dynamics.
\end{remark}

\begin{definition}[Affine homogeneous symmetry]
\label{def:affine-symmetry}
For \(c\in\R^3\), define
\[
   (S_cV)(y,\tau):=V(y-2c,\tau)-c,
   \qquad
   (S_cP)(y,\tau):=P(y-2c,\tau)+\frac12 c\cdot y .
\]
\end{definition}

\begin{lemma}[Affine renormalized Galilean symmetry]
\label{lem:affine-symmetry}
If \((V,P)\) solves \eqref{eq:centered-main}, then
\((S_cV,S_cP)\) also solves \eqref{eq:centered-main}.
\end{lemma}

\begin{proof}
Set \(U(y,\tau)=V(y-2c,\tau)-c\),
\(\Pi(y,\tau)=P(y-2c,\tau)+\frac12c\cdot y\), and \(z=y-2c\).  Then
\[
   \partial_\tau U=(\partial_\tau V)(z,\tau),
   \qquad \Delta U=(\Delta V)(z,\tau),
   \qquad \nabla U=(\nabla V)(z,\tau).
\]
The nonlinear term is
\[
   (U\cdot\nabla)U=(V-c)\cdot\nabla V
   =(V\cdot\nabla)V-c\cdot\nabla V,
\]
where all quantities on the right are evaluated at \((z,\tau)\).  The drift
term satisfies
\[
   \frac12 y\cdot\nabla U
   =\frac12(z+2c)\cdot\nabla V
   =\frac12z\cdot\nabla V+c\cdot\nabla V.
\]
The two \(c\cdot\nabla V\) terms cancel.  Finally,
\[
   \frac12U+\nabla\Pi
   =
   \frac12(V-c)+\nabla P+\frac12c
   =
   \frac12V+\nabla P.
\]
Substitution into \eqref{eq:centered-main} at the point \((z,\tau)\) gives the
equation for \((U,\Pi)\).  The divergence condition is unchanged by translation
and subtracting a constant vector.
\end{proof}

\begin{definition}[Asymptotic hull at infinity]
\label{def:asymptotic-hull}
Let \(V\) be a bounded smooth ancient solution of \eqref{eq:centered-main}.
Define
\[
   \calA_\infty(V)
   :=
   \left\{
      W:\ \exists a_n\in\R^3,\ |a_n|\to\infty,\ 
      T_{a_n}V\to W
      \text{ in } C^\infty_{\loc}(\R^3\times\R)
   \right\}.
\]
\end{definition}

\begin{theorem}[Compactness of the asymptotic hull]
\label{thm:asymptotic-hull-compactness}
Let \(V\) be a bounded smooth ancient solution of \eqref{eq:centered-main}.
Then \(\calA_\infty(V)\) is nonempty and compact in
\(C^\infty_{\loc}(\R^3\times\R)\).  Every
\(W\in\calA_\infty(V)\) is a bounded smooth ancient solution of
\eqref{eq:centered-main} with
\(\|W\|_{L^\infty}\le\|V\|_{L^\infty}\).
\end{theorem}

\begin{proof}
For any sequence \(a_n\) with \(|a_n|\to\infty\), the translated solutions
\(T_{a_n}V\) solve the same equation and have the same \(L^\infty\) bound by
\Cref{lem:covariant-translation-invariance}.  On every compact cylinder, the
coefficients of \eqref{eq:centered-main} are fixed and smooth.  Interior
parabolic estimates for the Stokes system with bounded smooth nonlinear data,
combined with the pressure equation
\(-\Delta P=\partial_i\partial_j(V_iV_j)\) after local gauges, give uniform
\(C^k\)-bounds on smaller compact cylinders for every \(k\).  A diagonal
Arzela--Ascoli argument gives a subsequential limit in
\(C^\infty_{\loc}\).  Passing to the limit in the equation gives another
bounded ancient solution.  The same compactness argument applied to an
arbitrary sequence of elements of \(\calA_\infty(V)\) proves compactness.
\end{proof}

\begin{definition}[Tail alternatives]
\label{def:R4-tail-strata}
Let \(\calC\) denote the set of constant vector fields.  Define
\[
   {\Cgen}_{\mathrm{tail\text{-}const}}
   :=
   \{V\in\Cgen:\calA_\infty(V)\subset\calC\},
\]
and
\[
   {\Cgen}_{\mathrm{tail\text{-}prof}}
   :=
   \{V\in\Cgen:\calA_\infty(V)\hbox{ contains a nonconstant profile}\}.
\]
\end{definition}

\begin{proposition}[Tail dichotomy and tail-zero normalization]
\label{prop:tail-dichotomy}
Every \(V\in\Cgen\) belongs to
\({\Cgen}_{\mathrm{tail\text{-}const}}\cup
{\Cgen}_{\mathrm{tail\text{-}prof}}\).  If
\(\calA_\infty(V)=\{C_\infty\}\) is a singleton constant hull, then
\[
   U:=S_{C_\infty}V
\]
solves \eqref{eq:centered-main} and satisfies
\[
   \calA_\infty(U)=\{0\}.
\]
\end{proposition}

\begin{proof}
The dichotomy is the definition.  Since \(T_a\) and \(S_c\) commute,
\[
   T_a(S_{C_\infty}V)=S_{C_\infty}(T_aV).
\]
Thus every far-field limit of \(S_{C_\infty}V\) is
\[
   S_{C_\infty}(C_\infty)=0.
\]
The equation is preserved by \Cref{lem:affine-symmetry}.
\end{proof}

\begin{definition}[Tail-minimal recurrent core]
\label{def:tail-minimal-core}
A nonempty compact set \(\calM\subset\calA_\infty(V)\) is a tail-minimal
recurrent core if it is invariant under all covariant translations \(T_a\),
invariant under all time translations
\[
   (\Theta_sW)(y,\tau):=W(y,\tau+s),
\]
and contains no proper nonempty compact subset with the same two invariance
properties.
\end{definition}

\begin{lemma}[Pressure gradients in invariant local averages]
\label{lem:recurrent-core-pressure-gauges}
Let \(\calM\) be a compact set of smooth bounded solutions of
\eqref{eq:centered-main}, compact in the locally smooth topology and invariant
under the covariant translations \(T_a\) and the time translations
\(\Theta_s\).  Let \(\mu\) be a Borel probability measure on \(\calM\)
invariant under both actions.  For each \(W\in\calM\), let \(P_W\) be any
smooth local pressure near \((0,0)\) satisfying \eqref{eq:centered-main}.  Put
\[
   u_i(W):=W_i(0,0),
   \qquad
   p_i(W):=\partial_iP_W(0,0),
   \qquad
   b_{ij}(W):=u_i(W)u_j(W).
\]
Let \(\Pi_0\) denote the orthogonal projection in
\(L^2(\calM,\mu)\) onto the closed subspace of functions invariant under all
spatial covariant translations \(T_a\), and set
\[
   u_i^0:=\Pi_0u_i,\qquad p_i^0:=\Pi_0p_i.
\]
Then the pressure contribution in the recurrent-core energy average is
gauge-independent and satisfies
\begin{equation}\label{eq:pressure-average-core}
   \int_{\calM}u_i(W)p_i(W)\,d\mu(W)
   =
   \int_{\calM}u_i^0(W)p_i^0(W)\,d\mu(W).
\end{equation}
Equivalently, the nonzero spatial-frequency part of the pressure gradient
does no work in invariant average:
\begin{equation}\label{eq:pressure-nonzero-modes-core}
   \int_{\calM}u_i(W)\bigl(p_i(W)-p_i^0(W)\bigr)\,d\mu(W)=0 .
\end{equation}
\end{lemma}

\begin{proof}
The additive pressure gauge is irrelevant because only \(\nabla P_W\) enters.
The pressure gradient is fixed by the velocity through the equation
\[
   \nabla P_W
   =
   \Delta W-\frac12(W+y\cdot\nabla W)
   -\partial_\tau W-(W\cdot\nabla)W .
\]
Thus \(W\mapsto p_i(W)\) is a continuous bounded local observable on
\(\calM\).  Paper II, \cite[Lemma~\emph{Local pressure reconstruction}]
{DuranPaperII}, supplies the local pressure representatives; the present
argument uses only their gradients, so the additive gauges never enter.

Let \(U_a\) be the unitary representation of spatial covariant translations
on \(L^2(\calM,\mu)\),
\[
   (U_aF)(W):=F(T_aW).
\]
Let \(D_1,D_2,D_3\) be its skew-adjoint infinitesimal generators.  For a smooth
local observable \(F\), \(D_kF\) is obtained by differentiating
\(F(T_{he_k}W)\) at \(h=0\).  In particular, at \(\tau=0\),
\[
   D_k u_i(W)=\partial_kW_i(0,0).
\]
Hence
\begin{equation}\label{eq:state-div-free}
   D_i u_i=0
   \qquad\hbox{in }L^2(\calM,\mu),
\end{equation}
because each \(W\) is divergence-free.

Taking the spatial divergence of \eqref{eq:centered-main} gives the local
pressure Poisson equation
\begin{equation}\label{eq:local-pressure-poisson-core}
   -\Delta_yP_W=\partial_i\partial_j(W_iW_j).
\end{equation}
Indeed, the divergence of \(\partial_\tau W\), \(\Delta W\), and \(W\) is
zero; the divergence of \(y\cdot\nabla W\) is
\(\partial_i(y_j\partial_jW_i)=\partial_jW_j+y_j\partial_j\partial_iW_i=0\);
and \((W\cdot\nabla)W_i=\partial_j(W_jW_i)\).  Applying the infinitesimal
translation generators at \((0,0)\) gives
\begin{equation}\label{eq:state-pressure-divergence}
   D_i p_i=-D_iD_jb_{ij}.
\end{equation}
The pressure gradient is curl-free, hence
\begin{equation}\label{eq:state-pressure-curl}
   D_kp_i=D_ip_k .
\end{equation}

Let
\[
   A:=-\sum_{k=1}^3D_k^2.
\]
By the spectral theorem for the commuting unitary group \(U_a\), \(A\) is
self-adjoint and nonnegative, and \(\Pi_0\) is exactly the spectral projection
of \(A\) at the eigenvalue \(0\).  Define the nonzero spatial part
\[
   q_i:=p_i-p_i^0=(I-\Pi_0)p_i .
\]
Equations \eqref{eq:state-pressure-divergence} and
\eqref{eq:state-pressure-curl} imply, on the nonzero spectral subspace,
\[
   q_k
   =
   D_kA^{-1}(I-\Pi_0)D_iD_jb_{ij}
\]
in the standard regularized sense.  Explicitly, with
\[
   q_{k,\varepsilon}
   :=
   D_k(\varepsilon+A)^{-1}(I-\Pi_0)D_iD_jb_{ij},
   \qquad \varepsilon>0,
\]
the joint spectral resolution of \(D_1,D_2,D_3\) gives
\[
   q_{k,\varepsilon}\longrightarrow q_k
   \quad\hbox{in }L^2(\calM,\mu)
   \quad\hbox{as }\varepsilon\downarrow0.
\]
The convergence is just the multiplier convergence
\[
   \frac{\xi_k\xi_i\xi_j}{\varepsilon+|\xi|^2}
      \longrightarrow
   \frac{\xi_k\xi_i\xi_j}{|\xi|^2}
   \qquad (\xi\ne0),
\]
together with \((I-\Pi_0)\), which removes the point \(\xi=0\).  The
curl-free relation \eqref{eq:state-pressure-curl} says that the nonzero
spatial spectral part of \(p\) is parallel to \(\xi\), and
\eqref{eq:state-pressure-divergence} fixes its longitudinal coefficient;
there is therefore no additional harmonic component.

For each \(\varepsilon>0\), skew-adjointness of \(D_k\) and
\eqref{eq:state-div-free} give
\begin{align*}
   \int_{\calM}u_k q_{k,\varepsilon}\,d\mu
   &=
   \int_{\calM}
      u_kD_k(\varepsilon+A)^{-1}(I-\Pi_0)D_iD_jb_{ij}\,d\mu \\
   &=
   -\int_{\calM}
      (D_ku_k)(\varepsilon+A)^{-1}(I-\Pi_0)D_iD_jb_{ij}\,d\mu
   =0 .
\end{align*}
Passing to the limit in \(L^2\) yields
\[
   \int_{\calM}u_k q_k\,d\mu=0,
\]
which is \eqref{eq:pressure-nonzero-modes-core}.  Since \(p_i^0\) lies in the
range of \(\Pi_0\) and \(u_i-u_i^0\) is orthogonal to that range,
\[
   \int_{\calM}(u_i-u_i^0)p_i^0\,d\mu=0 .
\]
Combining this with \eqref{eq:pressure-nonzero-modes-core} proves
\eqref{eq:pressure-average-core}.
\end{proof}

\begin{theorem}[Rigidity of recurrent profile tail cores]
\label{thm:recurrent-tail-core-rigidity}
Let \(\calM\) be a compact tail-minimal recurrent core for
\eqref{eq:centered-main} arising from the Paper I normalized Seregin theorem,
with the affine zero mode interpreted by the affine symmetry
\Cref{lem:affine-symmetry}.  Equivalently, after applying one fixed affine
normalization \(S_c\), the spatially homogeneous part of the core satisfies
the bounded mild centered formulation of Paper II
\cite[Definition~\emph{Bounded mild ancient solution}]{DuranPaperII}.  Then
every \(W\in\calM\) is spatially and temporally constant.
More precisely, there exists \(m\in\R^3\) such that
\[
   W(y,\tau)\equiv m
   \qquad
   \hbox{for every }W\in\calM.
\]
After applying the affine symmetry \(S_m\), the core is reduced to the zero
solution.
\end{theorem}

\begin{proof}
The group generated by covariant translations and time translations is
amenable on the compact invariant set \(\calM\).  Choose an invariant Borel
probability measure \(\mu\) with full support on a minimal component, and write
\[
   \langle F\rangle:=\int_{\calM}F(W)\,d\mu(W)
\]
for local observables.  Let \(P_W\) be a local pressure associated to \(W\).
As in \Cref{lem:recurrent-core-pressure-gauges}, set
\[
   u_i(W):=W_i(0,0),
   \qquad
   p_i(W):=\partial_iP_W(0,0),
\]
and let \(\Pi_0\) be the orthogonal projection onto observables invariant under
all spatial covariant translations.  Write
\[
   u_i^0:=\Pi_0u_i,
   \qquad
   p_i^0:=\Pi_0p_i .
\]
Invariance gives
\[
   \langle \partial_\tau F(W)(0,0)\rangle=0,
   \qquad
   \langle \partial_j F(W)(0,0)\rangle=0
\]
for every smooth local observable \(F\).  These identities follow first for
difference quotients generated by the group action and then by smoothness.

Evaluate \eqref{eq:centered-main} at \((0,0)\).  The drift term vanishes at
\(y=0\), and we get
\[
   \partial_\tau W(0,0)-\Delta W(0,0)+\frac12 W(0,0)
   +(W(0,0)\cdot\nabla)W(0,0)+\nabla P_W(0,0)=0 .
\]
Project this identity onto the spatially invariant subspace.  The spatial
derivative terms vanish under \(\Pi_0\): the Laplacian is
\(D_jD_j u_i\), and the nonlinear term is \(D_j(u_ju_i)\) because
\(D_ju_j=0\).  Since time translations commute with spatial translations,
\(\Pi_0\partial_\tau u_i=\partial_\tau u_i^0\).  Thus
\begin{equation}\label{eq:zero-mode-momentum-core}
   \partial_\tau u_i^0+\frac12u_i^0+p_i^0=0
   \qquad\hbox{in }L^2(\calM,\mu).
\end{equation}
Taking the inner product of \eqref{eq:zero-mode-momentum-core} with \(u_i^0\),
summing in \(i\), and using invariance under time translations,
\[
   \left\langle u_i^0\partial_\tau u_i^0\right\rangle
   =
   \frac12\left\langle\partial_\tau |u^0|^2\right\rangle
   =0,
\]
gives
\begin{equation}\label{eq:zero-mode-pressure-work-core}
   \left\langle u_i^0p_i^0\right\rangle
   =
   -\frac12\left\langle|u^0|^2\right\rangle .
\end{equation}

Dot the equation with \(W\), evaluate at \((0,0)\), and average.  The time
term is
\[
   \left\langle W\cdot\partial_\tau W\right\rangle
   =\frac12\left\langle\partial_\tau |W|^2\right\rangle=0 .
\]
The diffusion term gives
\[
   -W\cdot\Delta W
   =
   |\nabla W|^2-\divergence(W\cdot\nabla W),
\]
and therefore
\[
   -\left\langle W\cdot\Delta W\right\rangle
   =
   \left\langle|\nabla W|^2\right\rangle .
\]
The nonlinear term is a divergence:
\[
   W\cdot(W\cdot\nabla)W
   =
   \frac12W\cdot\nabla |W|^2
   =
   \frac12\divergence(|W|^2W),
\]
so its average is zero.  Decompose
\[
   p_i=p_i^0+(p_i-p_i^0).
\]
The nonzero spatial part of the pressure gradient does no work in invariant
average by \Cref{lem:recurrent-core-pressure-gauges}, and the remaining
spatial zero mode is evaluated by
\eqref{eq:zero-mode-pressure-work-core}.  Hence
\[
   \left\langle u_ip_i\right\rangle
   =
   \left\langle u_i^0p_i^0\right\rangle
   =
   -\frac12\left\langle|u^0|^2\right\rangle .
\]
Thus
\begin{equation}\label{eq:mean-energy-core}
   \left\langle|\nabla W|^2\right\rangle
   +\frac12\left\langle|W|^2\right\rangle
   -\frac12\left\langle|u^0|^2\right\rangle=0.
\end{equation}
Both terms are nonnegative.  The second nonnegative quantity is
\[
   \left\langle|u|^2\right\rangle-\left\langle|u^0|^2\right\rangle
   =
   \left\langle|u-u^0|^2\right\rangle,
\]
because \(\Pi_0\) is an orthogonal projection in \(L^2(\calM,\mu)\).
Therefore
\[
   \left\langle|\nabla W|^2\right\rangle=0,
   \qquad
   \left\langle|u-u^0|^2\right\rangle=0 .
\]
Since \(\mu\) has full support and the observables are continuous in the local
smooth topology, \(\nabla W(0,0)=0\) for every \(W\in\calM\).  Invariance
moves \((0,0)\) to any \((y,\tau)\), so \(\nabla W(y,\tau)=0\) everywhere.
Thus every element of \(\calM\) is spatially constant:
\[
   W(y,\tau)=a_W(\tau).
\]
It remains only to identify the allowed spatially homogeneous zero mode.  This
is not a pressure-gauge assertion.  By the Paper I normalization quoted in the
statement, there is one fixed affine normalization \(S_c\) for which the
spatially homogeneous representatives satisfy the bounded mild centered
formulation of Paper II.  For a spatially constant mild centered solution
\(B(\tau)\), the nonlinear Duhamel term in Paper II's mild formula vanishes
and the Ornstein--Uhlenbeck--Stokes semigroup of Paper II acts on constants by
\[
   S(\tau-\sigma)B(\sigma)=e^{-(\tau-\sigma)/2}B(\sigma).
\]
Thus
\[
   B(\tau)=e^{-(\tau-\sigma)/2}B(\sigma)
   \qquad(\sigma<\tau).
\]
Letting \(\sigma\to-\infty\) and using boundedness of the ancient trajectory
gives \(B(\tau)=0\) for every \(\tau\).  Undoing the fixed affine
normalization gives a single vector \(m=c\) and
\[
   W(y,\tau)\equiv m
   \qquad(W\in\calM).
\]
The affine symmetry \(S_m\) then maps the core to the zero solution.
\end{proof}

\section{Terminal concentration profiles and separated centers}
\label{sec:terminal-concentration}

The terminal part of the proof is local in parabolic cylinders.  It does not require a
global \(L^3\) bound, a global CKN budget, a finite critical-mass budget, or a
global nonlinear profile decomposition for the original residual profile.
Everything is formulated on fixed compact terminal cylinders after applying the
recentering sequence under consideration.

\begin{definition}[Parabolic recentering sequence]
\label{def:parabolic-recentering-sequence}
Let \((V_n,P_n)\) be a terminal extraction sequence on expanding domains
\(D_n\uparrow \R^3\times\R\), with local suitability and pressure-gauge
control.  A parabolic recentering sequence is a sequence
\[
   \mathfrak c=(z_n)_{n\ge1},\qquad z_n=(y_n,\tau_n)\in D_n .
\]
Two recentering sequences are equivalent if
\[
   \sup_n\left(|y_n-y_n'|+|\tau_n-\tau_n'|^{1/2}\right)<\infty,
\]
and asymptotically separated if the same parabolic distance tends to infinity.
The recentered fields associated to \(\mathfrak c\) are denoted
\[
   V_n^{\mathfrak c}(y,s):=V_n(y+y_n,s+\tau_n),
   \qquad
   P_n^{\mathfrak c}(y,s):=
   P_n(y+y_n,s+\tau_n)-a_n^{\mathfrak c}(s),
\]
where \(a_n^{\mathfrak c}\) is an admissible time-dependent local pressure
gauge.  All convergence assertions for a recentering sequence are made on fixed
compact sets in the variables \((y,s)\).
\end{definition}

\begin{definition}[Local CKN measure and concentrating recentering sequences]
\label{def:terminal-measures}
After choosing local pressure gauges \(a_n(\tau)\), define
\[
   d\mu_n
   :=
   \left(|V_n|^3+|P_n-a_n(\tau)|^{3/2}\right)\,dy\,d\tau,
   \qquad
   d\nu_n:=|\nabla V_n|^2\,dy\,d\tau .
\]
The local mass of a recentering sequence \(\mathfrak c=(z_n)\) is
\[
   m(\mathfrak c):=\limsup_{n\to\infty}\mu_n(Q_1(z_n)).
\]
For a fixed small-data threshold \(\eps_{\rm sd}>0\), the recentering sequence
is \(\eps_{\rm sd}\)-concentrating if \(m(\mathfrak c)\ge\eps_{\rm sd}\), and
locally vanishing at unit scale otherwise.  It has loss of local compactness if
\[
   \limsup_{n\to\infty}\nu_n(Q_2(z_n))=\infty
\]
or if the local pressure gauges fail to be uniformly controlled in
\(L^{3/2}(Q_2(z_n))\).
\end{definition}

\begin{lemma}[Paper I compactness excludes retained local compactness failure]
\label{lem:rough-stratum-discharge}
Let \((V_n,P_n)\) be a terminal sequence obtained from a Paper I Type I
Seregin extraction, or from a later terminal extraction of such a sequence.
If a recentering sequence \(\mathfrak c=(z_n)\) carries retained compact
velocity \(L^3\)-mass, then
after passing to a subsequence it has no loss of local compactness and its
recentered fields
converge locally in the Seregin topology to a suitable bounded terminal
profile.
\end{lemma}

\begin{proof}
Fix \(R<\infty\).  The recentered cylinders \(Q_{2R}(z_n)\) are compact
terminal cylinders in the variables of the chosen extraction.  Paper I gives
local suitability, local pressure-gauge compactness, and the inherited Type I
bound on every such compact cylinder.  Therefore, after choosing the pressure
gauges \(a_n^{\mathfrak c}\), the sequence
\((V_n^{\mathfrak c},P_n^{\mathfrak c})\) has the standard local bounds needed
for the Seregin compactness theorem on \(Q_R(0,0)\): strong local \(L^3\)
compactness for the velocity, weak local \(L^{3/2}\) compactness for the
pressure, and the local energy inequality in the limit.  A diagonal argument
over \(R=1,2,\dots\) gives a suitable bounded terminal profile.  The two
failures used in the definition of loss of local compactness are precisely the
negations of these local compactness and pressure-gauge conclusions, so a
retained recentering sequence arising from the Paper I Type I extraction theorem
cannot remain in the local-compactness-failure alternative.
\end{proof}

\begin{lemma}[Maximal separated families of concentrating recenterings]
\label{lem:maximal-active-family}
There exists a maximal separated family of \(\eps_{\rm sd}\)-concentrating
recentering sequences without loss of local compactness.  After this family is
extracted, every residual unit parabolic cylinder outside sufficiently large
tubes around the selected recentering sequences has CKN mass below
\(\eps_{\rm sd}\).
\end{lemma}

\begin{proof}
The existence of a maximal separated concentrating family follows by Zorn's
lemma on the partially ordered set of separated concentrating families.  If residual local
vanishing failed after extraction, there would be centers \(z_n'\) outside all
large selected tubes such that
\[
   \limsup_{n\to\infty}\mu_n(Q_1(z_n'))\ge\eps_{\rm sd}.
\]
Such a recentering sequence is \(\eps_{\rm sd}\)-concentrating and
asymptotically separated from every selected concentrating recentering sequence,
contradicting maximality.  Hence the residual after maximal concentration
extraction is locally vanishing at unit scale.
\end{proof}

\begin{lemma}[Separation of recentered cylinders]
\label{lem:separated-recentered-cylinders}
Let \(\mathfrak c_i=(z_{i,n})\) and \(\mathfrak c_j=(z_{j,n})\) be
asymptotically separated recentering sequences.  For every fixed \(R<\infty\), the
cylinder \(Q_R(z_{j,n})\), written in the coordinates of \(\mathfrak c_i\),
leaves every compact subset of the recentered variables.  Consequently
no positive compact velocity \(L^3\)-mass assigned to \(\mathfrak c_j\) is visible on any
fixed compact cylinder of the \(\mathfrak c_i\)-frame.
\end{lemma}

\begin{proof}
In the \(\mathfrak c_i\)-coordinates the center of \(Q_R(z_{j,n})\) is
\[
   \bigl(y_{j,n}-y_{i,n},\,\tau_{j,n}-\tau_{i,n}\bigr).
\]
The parabolic distance of this point from the origin tends to infinity by
asymptotic separation.  Hence, for every fixed compact terminal cylinder
\(K\), one has
\[
   K\cap\bigl(Q_R(z_{j,n})-z_{i,n}\bigr)=\varnothing
\]
for all sufficiently large \(n\).  This is a purely local statement about the
relative positions of the two recentering sequences.  It does not estimate the solution on
all of space and does not use a global norm.
\end{proof}

\begin{theorem}[Terminal local concentration alternatives]
\label{thm:terminal-exhaustion-main}
Every terminal concentration sequence belongs to exactly one of the following
alternatives:
\[
   \mathcal S_{\rm one},\quad
   \mathcal S_{\rm locvan},\quad
   \mathcal S_{\rm ext},\quad
   \mathcal S_{\rm diff},\quad
   \mathcal S_{\rm noncomp},\quad
   \mathcal S_{\rm sep}.
\]
Here \(\mathcal S_{\rm noncomp}\) consists of recentering sequences with loss of
local compactness; \(\mathcal S_{\rm one}\) and \(\mathcal S_{\rm sep}\) are
the bounded concentrating alternatives with respectively one, or at least two,
bounded separated equivalence classes; \(\mathcal S_{\rm ext}\) consists of
concentrating recentering sequences escaping the terminal frame; \(\mathcal
S_{\rm diff}\) consists of diffuse exterior concentration, namely critical mass
persisting on expanding regions while every fixed unit cylinder is locally
vanishing; and \(\mathcal S_{\rm locvan}\) is local CKN vanishing after removal
of the selected concentrating recentering sequences.
\end{theorem}

\begin{proof}
Let \(\mathfrak c\) be a terminal concentration sequence.  If it has loss of
local compactness, it belongs to \(\mathcal S_{\rm noncomp}\).  Assume it has no
loss of local compactness.  If it is \(\eps_{\rm sd}\)-concentrating, maximal
concentration extraction puts it into the selected concentrating family up to
equivalence.
If its representative remains bounded in the terminal frame and it is the
unique bounded concentrating class, it gives the single-profile alternative.  If
more than one bounded concentrating class is present, it gives a finite separated
profile family.  If it escapes every bounded terminal frame, it is the exterior
escape alternative.  This classifies concentrating sequences without loss of
local compactness.

It remains to classify residual sequences which are locally vanishing at unit
scale and have no loss of local compactness.  By
\Cref{lem:maximal-active-family}, the residual after concentration extraction
has local CKN vanishing on every fixed compact terminal cylinder.  If no
critical mass remains on expanding regions, this is the local-vanishing
alternative \(\mathcal S_{\rm locvan}\).  If critical mass persists only on
expanding regions while every fixed unit cylinder remains below
\(\eps_{\rm sd}\), this is the diffuse exterior concentration alternative
\(\mathcal S_{\rm diff}\).  This is a terminal alternative, not an a priori
global estimate on the original solution.  The alternatives are mutually
exclusive by construction and exhaustive by the preceding cases.
\end{proof}

\begin{corollary}[Persistent local velocity mass forces retained concentration]
\label{cor:persistent-forces-active}
If a terminal sequence satisfies
\[
   \liminf_{n\to\infty}
   \iint_{Q_1(0,0)} |V_n|^3\,dy\,d\tau\ge\eta_0>0
\]
with \(\eta_0\ge\eps_{\rm sd}\), then the terminal decomposition contains a
local concentration profile unless the distinguished recentering sequence has
loss of local compactness.  In the generic residual after local vanishing and
local compactness failure are excluded, the persistent compact velocity mass is
therefore carried by either one retained concentration profile or a finite
separated family of retained concentration profiles.
\end{corollary}

\begin{proof}
The local CKN measure \(\mu_n\) dominates the velocity contribution.  Hence the
origin recentering sequence has \(\mu\)-mass at least \(\eta_0\), and is
\(\eps_{\rm sd}\)-concentrating unless it has loss of local compactness.  If it
has no loss of local compactness and is bounded in the terminal frame, it lies
in the single-profile or separated-profile alternative of
\Cref{thm:terminal-exhaustion-main}.
\end{proof}

\begin{lemma}[Local vanishing and exterior alternatives do not carry compact concentration]
\label{lem:radiative-stratum-discharge}
Let a terminal extraction inherit the positive compact velocity \(L^3\)-mass
supplied by Paper 0 and Paper I.  Then the retained compact velocity mass
cannot be carried solely by local vanishing, exterior escape, or diffuse
exterior concentration.
\end{lemma}

\begin{proof}
The assertion is local.  By Paper 0, failure of local vanishing at a singular
terminal point gives a positive local velocity concentration sequence.  Paper I
carries this through the Type I Seregin extraction as the compact velocity lower
bound \eqref{eq:positive-velocity-main}.  Cover the fixed compact cylinder
\(K\) in \eqref{eq:positive-velocity-main} by finitely many unit terminal
cylinders.  At least one cylinder in the cover carries a positive velocity
lower bound, say \(\eta_K>0\).  Since \(\mu_n\) dominates
\(|V_n|^3\,dy\,d\tau\), choosing \(\eps_{\rm sd}\le\eta_K\) makes that unit
cylinder \(\eps_{\rm sd}\)-concentrating.  Thus the distinguished compact
concentration cannot be locally vanishing.

Local vanishing after removal of concentrating recentering sequences means
that the residual state has local CKN mass below \(\eps_{\rm sd}\) on each fixed
compact terminal cylinder; it therefore cannot contain the fixed compact
velocity lower bound.  An exterior escaping concentration is visible only after
its own recentering; relative to the distinguished compact terminal cylinder it
escapes every bounded window.  By
\Cref{lem:separated-recentered-cylinders}, it cannot be the sole location of a
lower bound that remains in the distinguished compact window.  Finally, diffuse
exterior concentration is defined by persistence through expanding regions
while all fixed unit cylinders are locally vanishing.  If such a sequence
contained the compact Paper 0/Paper I velocity lower bound, the positive
concentration lemma of Paper 0 would select a unit terminal cylinder with mass
at least the retained threshold, contradicting local vanishing.  Thus none of
these alternatives can be the unique location of the retained compact velocity
mass.
\end{proof}

\begin{theorem}[Local exclusion of non-compact and non-concentrating terminal alternatives]
\label{thm:terminal-nonactive-discharge}
For terminal sequences generated by the Paper I Type I Seregin extraction
theorem, loss of local compactness, local vanishing, exterior escape, and
diffuse exterior concentration are not possible sole locations of the retained
compact velocity mass.
\end{theorem}

\begin{proof}
Retained recentering sequences with loss of local compactness are removed by
\Cref{lem:rough-stratum-discharge}.  Local vanishing, exterior escape, and
diffuse exterior concentration are removed by
\Cref{lem:radiative-stratum-discharge}.  The proof uses only the local positive
concentration theorem of Paper 0, the local Type I Seregin compactness and
pressure gauges of Paper I, and fixed compact terminal cylinders.
\end{proof}

\begin{definition}[Finite separated family of retained concentration profiles]
\label{def:finite-separated-profile-family}
A finite separated family of retained concentration profiles is a finite
family of parabolic recentering sequences
\(g_{1,n},\dots,g_{J,n}\), \(J\ge2\), such that:
\begin{enumerate}[label=(\roman*)]
   \item the recentering sequences are mutually separated;
   \item each recentered sequence converges locally smoothly for velocity and
   locally in \(L^{3/2}\) for pressure modulo gauges to a terminal profile
   \((U_j,\Pi_j)\);
   \item each limiting profile has retained local concentration:
\[
   \calE_1(U_j,\Pi_j;0)\ge\eps_*
   \qquad (j=1,\dots,J).
\]
\end{enumerate}
This is a finite local family, not a statement that all concentration profiles
in the whole space have been counted.  The separated-profile proof below treats
the possible continuation of escaping recenterings by local concentration
alternatives, not by a global counting estimate.
\end{definition}

\begin{proposition}[Finite separated families reduce to single-profile alternatives]
\label{prop:multi-no-new-obstruction}
Assume the single-profile retained concentration closure statement: no
single retained concentration profile remains in the refined residual decomposition after the already
closed local classes are removed.  Then a finite separated family is not
a new refined residual alternative.
\end{proposition}

\begin{proof}
Let \(A=\{\mathfrak s_1,\dots,\mathfrak s_J\}\) be a finite separated
concentrating
family, represented by mutually separated parabolic recentering sequences.
Recenter the terminal sequence at the sequence defining \(\mathfrak s_j\).
In that frame all other concentration profiles leave every compact terminal
cylinder, so the local terminal limit is a single retained concentration
profile.  By the single-profile closure statement, this component either belongs to
the already closed local classes, has local vanishing, or has loss of local
compactness.  Since the original family has retained concentration and local
compactness, the latter two alternatives are incompatible with the hypotheses.
Applying this argument to each \(j\) shows that the finite separated family
consists only of already closed local profiles, and hence does not remain in the
refined residual decomposition.
\end{proof}

\section{Separated concentration profiles, indecomposability, and sequence-\texorpdfstring{\(L^3\)}{L3}}
\label{sec:separated-profiles-atomic}

We now close the terminal alternatives carrying retained local concentration.
The proof uses no global critical-mass budget.  Instead, finite separated
concentration families are treated
through local recentering and compactness.  Iterating retained tail limits from
such a finite family has only
three local outcomes: an indecomposable retained concentration profile, a
finite recurrent tail loop, or escape into one of the non-concentrating or
non-compact alternatives already excluded by
\Cref{thm:terminal-nonactive-discharge}.

\begin{definition}[Tail limits and terminal indecomposability]
\label{def:atomic-terminal-profile}
Let \((U,\Pi)\) be a terminal profile with retained local concentration.  A
retained concentrating tail limit is a terminal limit obtained from recentering
sequences escaping to infinity relative to \(U\), with retained local
concentration at the origin.  Diffuse exterior concentration means that there
are \(\tau_n\to-\infty\), \(R_n\to\infty\), and \(\eta>0\), such that
\[
   \int_{|y|>R_n}|U(y,\tau_n)|^3\,dy\ge\eta,
\]
while every unit-scale exterior recentering sequence remains below the
retained-concentration threshold, unless that recentering sequence has loss of
local compactness.  An escaping non-compact recentering sequence is a tail
recentering sequence for which the local compactness, suitability, or pressure
control needed to extract a terminal profile fails.  A retained concentration
profile is terminally indecomposable if it has no retained concentrating tail
limit, no diffuse exterior concentration, and no escaping non-compact
recentering sequence.
\end{definition}

\begin{lemma}[One retained tail limit gives two separated concentration profiles]
\label{lem:descendant-to-separated-family}
Let \(U\) be a retained concentration terminal profile.  If \(U\) has a
retained concentrating tail limit, then the corresponding terminal extraction contains a
finite separated family of retained concentration profiles.
\end{lemma}

\begin{proof}
The origin recentering sequence carries retained concentration:
\[
   \calE_1(U,\Pi;0)\ge\eps_*.
\]
Let \(W\) be a retained concentrating tail limit.  By definition there are
relative tail recentering sequences \(h_n\to\infty\) such that \(h_n\cdot U\to W\) in the
terminal topology and
\[
   \calE_1(W,\Pi_W;0)\ge\eps_*.
\]
By local convergence of velocity and pressure after gauges,
\[
   \calE_1(h_n\cdot U,\Pi_{h_n};0)\ge\frac12\eps_*
\]
for all sufficiently large \(n\).  Thus, for large \(n\), the origin
recentering sequence and the sequence \(h_n\) are separated and have retained
local concentration.
They form a finite separated two-profile concentration family.
\end{proof}

\begin{proposition}[No finite separated family and local exclusions imply indecomposability]
\label{prop:no-multi-implies-atomic}
Every single retained concentration terminal profile arising from a refined
residual candidate is terminally indecomposable.
\end{proposition}

\begin{proof}
Let \(U\) be a single retained concentration terminal profile.  If \(U\) had a
retained concentrating tail limit, \Cref{lem:descendant-to-separated-family} would
produce a finite two-profile separated family with retained concentration, contradicting
\Cref{thm:no-finite-separated-profile-family}.  If \(U\) had diffuse exterior
concentration or an escaping non-compact recentering sequence, the corresponding
terminal extraction would belong to one of the terminal alternatives excluded
by \Cref{thm:terminal-nonactive-discharge}.  Hence no retained concentrating
tail limit, diffuse exterior concentration, or escaping non-compact recentering
sequence exists, and \(U\) is terminally indecomposable.
\end{proof}

\begin{lemma}[Indecomposable sequence-\(L^3\) completeness]
\label{lem:atomic-sequence-L3}
Let \((U,\Pi)\) be a bounded smooth ancient solution of
\eqref{eq:centered-main}.  Suppose \(U\) has retained local concentration and
is terminally indecomposable in the sense of
\Cref{def:atomic-terminal-profile}.  Then there exists a sequence
\(\tau_k\to-\infty\) such that
\[
   \sup_k\|U(\cdot,\tau_k)\|_{L^3(\R^3)}<\infty .
\]
\end{lemma}

\begin{proof}
We prove the contrapositive.  Define
\[
   A(\tau):=\int_{\R^3}|U(y,\tau)|^3\,dy\in[0,\infty].
\]
If the conclusion fails, then
\[
   \liminf_{\tau\to-\infty}A(\tau)=\infty.
\]
Choose \(\tau_n\to-\infty\) such that \(A(\tau_n)\to\infty\), allowing the
value \(+\infty\).  Since \(U\) is bounded by some \(M\), for each finite
\(R\),
\[
   \int_{|y|\le R}|U(y,\tau_n)|^3\,dy\le M^3|B_R|.
\]
Choose \(R_n\to\infty\) so slowly that, after discarding finitely many \(n\),
\[
   \int_{|y|>R_n}|U(y,\tau_n)|^3\,dy\ge1.
\]
If \(A(\tau_n)=+\infty\), take \(R_n=n\); otherwise choose \(R_n\) with
\(M^3|B_{R_n}|\le A(\tau_n)/2\).

Set
\[
   m_n:=\sup_{|x|>R_n}\int_{B_1(x)}|U(y,\tau_n)|^3\,dy .
\]
There are two alternatives.  First suppose
\(\limsup_{n\to\infty}m_n\ge\eps_*\).  Then, passing to a subsequence, there
are \(x_n\) with \(|x_n|>R_n\to\infty\) and
\[
   \int_{B_1(x_n)}|U(y,\tau_n)|^3\,dy\ge\eps_* .
\]
Use the covariant tail recentering
\[
   U_n(y,s):=U(y+e^{s/2}x_n,\tau_n+s).
\]
At \(s=0\),
\[
   \int_{B_1(0)}|U_n(y,0)|^3\,dy\ge\eps_* .
\]
If this recentering sequence has loss of local compactness, this contradicts
terminal indecomposability.  If it has local compactness, the compactness
theorem gives a subsequential smooth
limit \(W\), and
\[
   \int_{B_1(0)}|W(y,0)|^3\,dy\ge\eps_*,
\]
so \(W\) is a retained concentrating tail limit, again contradicting terminal
indecomposability.

The remaining alternative is that \(m_n<\eps_*\) for all sufficiently large
\(n\).  Together with the exterior lower bound, this says that at times
\(\tau_n\) there is at least one unit of \(L^3\)-mass outside \(B_{R_n}\), but
no unit-scale exterior recentering sequence carries retained local
concentration.  If some such exterior recentering sequence has loss of local
compactness, terminal indecomposability is contradicted.  If no such recentering
sequence has loss of local compactness, this is diffuse exterior concentration
with \(\eta=1\), again contradicting terminal indecomposability.  Both
alternatives are impossible, so the desired sequence exists.
\end{proof}

\begin{lemma}[Mildness inheritance for terminal profiles]
\label{hyp:mildness-inheritance-main}
If \(U\) is a terminal profile with retained local concentration, then the physical pullback
\[
   u(x,t)=(-t)^{-1/2}
   U\!\left(\frac{x}{\sqrt{-t}},-\log(-t)\right),
   \qquad t<0,
\]
is a mild ancient Navier--Stokes solution after any finite time shift.
\end{lemma}

\begin{proof}
Fix \(T<0\).  On \(\R^3\times(-\infty,T]\), the pullback is smooth and bounded
by \((-T)^{-1/2}\|U\|_{L^\infty}\).  It solves the physical Navier--Stokes
equations by the direct self-similar change of variables.  For
\(-\infty<s<t\le T\), applying the Duhamel formula to the smooth bounded
solution on \([s,t]\), with the Leray projection interpreted in the same
distributional-duality sense as the bounded mild formulation of Paper II
\cite{DuranPaperII}, gives
\[
   u(t)=e^{(t-s)\Delta}u(s)
   -\int_s^t e^{(t-\sigma)\Delta}
      \mathbb P\nabla\cdot(u\otimes u)(\sigma)\,d\sigma .
\]
After translating the terminal time \(T\) to \(0\), this is exactly the mild
ancient formulation on every finite forward interval.  Since \(T<0\) was
arbitrary, the pullback is mild after any finite time shift.
\end{proof}

\begin{theorem}[Albritton--Barker endpoint Liouville theorem]
\label{thm:AB-main}
Let \(u\) be a mild ancient solution of the three-dimensional Navier--Stokes
equations.  If there exists \(t_k\downarrow-\infty\) such that
\[
   \sup_k\|u(\cdot,t_k)\|_{L^3(\R^3)}<\infty,
\]
then \(u\equiv0\).
\end{theorem}

\begin{proof}
This is the endpoint ancient Liouville theorem in the sequence-\(L^3\) form
proved by Albritton and Barker \cite{AlbrittonBarker2019}.  Paper I records the
same peer-reviewed input among its auxiliary Liouville results
\cite{DuranPaperI}.
\end{proof}

\begin{lemma}[No retained indecomposable concentration profile]
\label{lem:no-atomic-active}
No retained terminally indecomposable concentration profile exists.
\end{lemma}

\begin{proof}
Let \(U\) have retained local concentration and be terminally indecomposable.  By
\Cref{lem:atomic-sequence-L3}, choose \(\tau_k\to-\infty\) such that
\[
   \sup_k\|U(\cdot,\tau_k)\|_{L^3}<\infty .
\]
Let \(u\) be the physical pullback and set \(t_k=-e^{-\tau_k}\).  Since
\(\tau_k\to-\infty\), the sequence satisfies \(t_k\downarrow-\infty\), and
critical scaling gives
\[
   \|u(\cdot,t_k)\|_{L^3(\R^3)}
   =
   \|U(\cdot,\tau_k)\|_{L^3(\R^3)}.
\]
By \Cref{hyp:mildness-inheritance-main}, \(u\) is mild ancient; hence
\Cref{thm:AB-main} gives \(u\equiv0\), and therefore \(U\equiv0\).  This
contradicts retained local concentration.
\end{proof}

\begin{lemma}[Diagonal lifting of retained concentrating tail limits]
\label{lem:diagonal-lifting-descendants}
Let \(A=\{\mathfrak s_1,\dots,\mathfrak s_J\}\), \(J\ge2\), be a finite
retained concentration separated family for a refined residual candidate, represented by
separated recentering sequences \(g_{j,n}=(y_{j,n},\tau_{j,n})\).  Let
\((U_j,\Pi_j)\) be the terminal profile obtained in the \(j\)-th recentering.  If
\((U_j,\Pi_j)\) has a retained concentrating tail limit \((W,\Pi_W)\), then
there is a parabolic recentering sequence \(\mathfrak d=(d_n)\) for the original terminal
sequence such that:
\begin{enumerate}[label=(\roman*)]
   \item \(\mathfrak d\) is asymptotically separated from \(g_j\);
   \item the recentered original sequence at \(\mathfrak d\) converges locally
   to \((W,\Pi_W)\), after admissible pressure gauges;
   \item \(\mathfrak d\) has retained local concentration.
\end{enumerate}
\end{lemma}

\begin{proof}
By definition of retained concentrating tail limit, there are relative tail
recenterings \(\zeta_m=(\xi_m,\sigma_m)\), with
\[
   |\xi_m|+|\sigma_m|^{1/2}\to\infty,
\]
such that the fields obtained from \((U_j,\Pi_j)\) by recentering at
\(\zeta_m\) converge locally in the terminal topology to \((W,\Pi_W)\), and
\[
   \calE_1(W,\Pi_W;0)\ge\eps_* .
\]
Fix \(m\).  Since \(V_n^{g_j}\to U_j\) locally smoothly and
\((P_n^{g_j})\to\Pi_j\) locally in \(L^{3/2}\) modulo gauges, the convergence
also holds on the fixed compact cylinder \(Q_m(\zeta_m)\).  Therefore one can
choose \(n(m)\) so large that the recentered original fields
\[
   V_{n(m)}(\,\cdot+y_{j,n(m)}+\xi_m,\,
             \cdot+\tau_{j,n(m)}+\sigma_m)
\]
and the corresponding pressure gauges are within \(m^{-1}\) of the
\(\zeta_m\)-recentered fields of \((U_j,\Pi_j)\) on \(Q_m(0,0)\).

Set
\[
   d_m:=(y_{j,n(m)}+\xi_m,\tau_{j,n(m)}+\sigma_m).
\]
Relabeling the subsequence \(n(m)\) as the terminal index gives a recentering sequence
\(\mathfrak d=(d_m)\) for the original terminal sequence.  Since
parabolic distance from \(d_m\) to \(g_{j,n(m)}\) is exactly
\(|\xi_m|+|\sigma_m|^{1/2}\), the sequence \(\mathfrak d\) is separated from
the \(j\)-th recentering sequence.  The diagonal choice gives local convergence of the
original sequence at \(\mathfrak d\) to \((W,\Pi_W)\).  Finally, the local
CKN quantity is continuous under the strong local velocity convergence and
weak lower semicontinuity of the pressure term after the fixed gauges, so
\[
   \liminf_{m\to\infty}\mu_{n(m)}(Q_1(d_m))
   \ge \calE_1(W,\Pi_W;0)\ge\eps_* .
\]
Thus \(\mathfrak d\) has retained local concentration.
\end{proof}

\begin{lemma}[Escaping recenterings in a separated family]
\label{lem:finite-successor-alternatives}
Let \(A=\{\mathfrak s_1,\dots,\mathfrak s_J\}\), \(J\ge2\), be a finite
retained concentration separated family for a refined residual candidate.  For each component
\(\mathfrak s_j\), exactly one of the following alternatives holds:
\begin{enumerate}[label=(\roman*)]
   \item the profile \((U_j,\Pi_j)\) is terminally indecomposable;
   \item \((U_j,\Pi_j)\) has diffuse exterior concentration or an escaping
   recentering sequence with loss of local compactness;
   \item \((U_j,\Pi_j)\) has a retained concentrating tail limit which lifts to a
   recentering sequence of the original terminal sequence with retained local
   concentration.
\end{enumerate}
\end{lemma}

\begin{proof}
This is the definition of terminal indecomposability plus the diagonal lifting
lemma.  If \((U_j,\Pi_j)\) is not terminally indecomposable, then at least one
of the three excluded alternatives in \Cref{def:atomic-terminal-profile} occurs:
a retained concentrating tail limit, diffuse exterior concentration, or an
escaping recentering sequence with loss of local compactness.  The diffuse
exterior and local-compactness-failure cases give (ii).  In the retained
concentrating tail-limit case, \Cref{lem:diagonal-lifting-descendants} lifts the
limit to a recentering sequence of the original terminal sequence with retained
local concentration, giving (iii).  The three alternatives are disjoint by
definition.
\end{proof}

\begin{lemma}[No diffuse or non-compact escaping recentering from a retained family]
\label{lem:no-nonactive-successor}
In the setting of \Cref{lem:finite-successor-alternatives}, alternative
\emph{(ii)} cannot occur.
\end{lemma}

\begin{proof}
Suppose that \((U_j,\Pi_j)\) has diffuse exterior concentration or an escaping
recentering sequence with loss of local compactness.  By the
same diagonal procedure used in
\Cref{lem:diagonal-lifting-descendants}, this alternative is realized by a
terminal extraction of the original refined residual candidate, after passing to
a subsequence and choosing the local pressure gauges supplied by Paper I.  Loss
of local compactness contradicts \Cref{lem:rough-stratum-discharge}.  Diffuse
exterior concentration places the retained compact velocity mass in the diffuse
exterior alternative, contradicting \Cref{lem:radiative-stratum-discharge}.
Therefore alternative (ii) cannot occur for a refined residual family with
retained local concentration.
\end{proof}

\begin{lemma}[Finite separated families produce indecomposable profiles or recurrent loops]
\label{lem:finite-graph-cycle}
Let \(A=\{\mathfrak s_1,\dots,\mathfrak s_J\}\), \(J\ge2\), be a finite
retained concentration separated family for a refined residual candidate.  Then either one of
the terminal profiles \((U_j,\Pi_j)\) is terminally indecomposable, or the
family generates a finite recurrent loop of retained concentrating tail limits.
\end{lemma}

\begin{proof}
Assume no \((U_j,\Pi_j)\) is terminally indecomposable.  By
\Cref{lem:finite-successor-alternatives,lem:no-nonactive-successor}, each
\((U_j,\Pi_j)\) has a retained concentrating tail limit which lifts to a
recentering sequence of the original terminal sequence with retained local
concentration.

If a lifted concentrating recentering sequence is not equivalent to one of the
sequences already in the finite family, enlarge the family by adding this new
sequence.  This is still a finite retained concentration separated family because only
one separated recentering sequence with retained local concentration has been added.  Repeat the
construction.  At each finite stage there
are two possibilities.  Either an indecomposable profile appears, proving the
first alternative, or every existing decomposable profile has a retained
concentrating tail limit.

If the process returns to a previously occurring concentrating recentering
sequence, the directed recentering relation contains a finite loop
\[
   \mathfrak s_{j_0}\to\mathfrak s_{j_1}\to\cdots
   \to\mathfrak s_{j_q}=\mathfrak s_{j_0}.
\]
This is the desired finite recurrent loop of retained concentrating tail limits.

It remains to explain why the construction cannot avoid both alternatives by
producing concentrating recentering sequences escaping every finite stage.  Such
an escaping sequence of lifted concentrating recenterings is, by
\Cref{thm:terminal-exhaustion-main}, an exterior escaping concentration
relative to each previously fixed finite family unless local compactness fails
or the mass becomes diffuse below the unit concentration threshold.  The latter
two cases are loss of local compactness and diffuse exterior concentration,
already excluded by \Cref{lem:no-nonactive-successor}.  The exterior case gives
a retained compact CKN concentration outside every fixed finite family.
Recenter at that exterior concentration.  In that frame, all previously chosen
finite family components leave
every compact terminal cylinder by \Cref{lem:separated-recentered-cylinders}; hence
the retained compact density is located solely in the exterior escaping
alternative relative to the preceding finite family.  This contradicts the exterior case of
\Cref{lem:radiative-stratum-discharge}.  Therefore indefinite escape without a
return is impossible, and a finite recurrent loop must occur unless an
indecomposable profile has already occurred.
\end{proof}

\begin{lemma}[Finite recurrent loops of retained tail limits are impossible]
\label{lem:no-finite-recurrent-cycle}
No refined residual terminal extraction contains a finite recurrent loop of retained
concentrating tail limits.
\end{lemma}

\begin{proof}
Let
\[
   \mathfrak s_{j_0}\to\mathfrak s_{j_1}\to\cdots
   \to\mathfrak s_{j_q}=\mathfrak s_{j_0}
\]
be such a loop, and write \((U_\ell,\Pi_\ell)\) for the corresponding retained
concentration terminal profiles.  Each arrow means that \(U_{\ell+1}\) is a
retained concentrating tail limit of \(U_\ell\).  Thus, for each \(\ell\), there are escaping
tail recentering sequences whose recentered fields converge locally to \(U_{\ell+1}\).

Let \(\calK\) be the closure, in the local smooth velocity topology and local
\(L^{3/2}\) pressure topology modulo gauges, of all profiles obtained from the
finite set \(\{U_0,\dots,U_{q-1}\}\) by time translations and covariant
terminal translations.  The boundedness inherited from the Type I Seregin
extraction theorem and the local pressure compactness of Paper I give compactness of
\(\calK\) on every fixed terminal cylinder; a diagonal argument gives compactness
of \(\calK\).  The set is invariant under the same time and covariant
translations by construction.  The recurrent-loop relations imply that each \(U_\ell\)
returns to \(\calK\) through tail limits, so the closed invariant subset of
\(\calK\) generated by the loop is nonempty, compact, invariant, and contains
the retained concentration profiles.

Choose a compact invariant subset \(\calM\) of this loop-generated set which
is minimal among compact invariant subsets still containing the loop profiles.
The invariant mean argument in
\Cref{thm:recurrent-tail-core-rigidity} applies to \(\calM\): all ingredients
used there are local on compact cylinders, and the pressure term is justified
by the invariant local pressure-gauge lemma
\Cref{lem:recurrent-core-pressure-gauges}.  Hence every element of \(\calM\),
in particular every loop profile \(U_\ell\), is spatially and temporally
constant.  Applying the affine homogeneous symmetry
\Cref{lem:affine-symmetry} removes this constant and sends each loop profile to
the zero profile.  This contradicts retained local concentration,
\[
   \calE_1(U_\ell,\Pi_\ell;0)\ge\eps_* .
\]
Therefore no finite recurrent loop of retained concentrating tail limits exists.
\end{proof}

\begin{theorem}[Finite separated retained concentration profile families are impossible]
\label{thm:no-finite-separated-profile-family}
No refined residual candidate admits a finite separated family of retained
concentration profiles.
\end{theorem}

\begin{proof}
Suppose, for contradiction, that a refined residual candidate admits a finite
retained concentration separated family \(A=\{\mathfrak s_1,\dots,\mathfrak s_J\}\),
\(J\ge2\).  By \Cref{lem:finite-graph-cycle}, either some component of the
family is terminally indecomposable
or the family generates a finite recurrent loop of retained concentrating tail
limits.  The first
alternative contradicts \Cref{lem:no-atomic-active}.  The second contradicts
\Cref{lem:no-finite-recurrent-cycle}.  Hence the finite retained concentration
separated family cannot exist.
\end{proof}

\begin{theorem}[Terminal local concentration theorem]
\label{thm:terminal-local-stratification-package}
Every bounded centered Type I Seregin profile with retained local concentration outside the
previously closed lower classes admits the local terminal concentration decomposition of
\Cref{thm:terminal-exhaustion-main}; loss of local compactness, local
vanishing, exterior escape, and diffuse exterior concentration are excluded by
\Cref{thm:terminal-nonactive-discharge}; finite separated retained
concentration profile families are excluded by
\Cref{thm:no-finite-separated-profile-family}; every remaining single retained
concentration terminal profile is terminally indecomposable by
\Cref{prop:no-multi-implies-atomic}; retained indecomposable profiles have a
backward sequence of finite critical norm by \Cref{lem:atomic-sequence-L3}; and
retained concentration terminal profiles inherit mildness by
\Cref{hyp:mildness-inheritance-main}.
\end{theorem}

\begin{proof}
This is the assembly of the local results proved in
\Cref{sec:terminal-concentration,sec:separated-profiles-atomic}.  The local concentration decomposition is
\Cref{thm:terminal-exhaustion-main}.  The local-vanishing, exterior-escape,
diffuse-exterior-concentration, and local-compactness-failure alternatives are
excluded by \Cref{thm:terminal-nonactive-discharge}.  The finite separated
concentration profile-family alternative is empty by
\Cref{thm:no-finite-separated-profile-family}.  Therefore any retained
concentration terminal profile that survives the ordered lower class removals is
single.  By \Cref{prop:no-multi-implies-atomic} it is terminally
indecomposable, and then
\Cref{lem:atomic-sequence-L3,hyp:mildness-inheritance-main} give the
sequence-\(L^3\) and mildness conclusions.  All estimates used in this assembly
are localized to fixed terminal cylinders, except for the final sequence
\(L^3\) conclusion, which is derived by contradiction from terminal
indecomposability and is not assumed as a global bound.
\end{proof}

\section{Closure of the generic terminal residual}
\label{sec:R4-closure}

\begin{definition}[Generic residual profile]
\label{def:R4-profile-main}
A generic residual profile is a smooth bounded normalized Seregin ancient solution of
\eqref{eq:centered-main} which carries persistent compact velocity
\(L^3\)-mass and
belongs to none of the previously closed alternatives, none of
\(\Cax,\Crot,\Cstat\), and none of the structured or tight
classes removed earlier.
\end{definition}

\begin{theorem}[Closure of the generic residual class by local concentration analysis]
\label{thm:R4-final-closure}
One has
\[
   \Cgen=\varnothing .
\]
\end{theorem}

\begin{proof}
Suppose, for contradiction, that \(V\in\Cgen\).  By the tail geometry
developed above, constant tails may be removed by the exact affine symmetry
and nonconstant recurrent profile tail cores collapse to affine constants.
Thus the remaining representative may be taken to be tail-normalized without
leaving the centered equation or the residual decomposition.  The inherited compact
velocity \(L^3\)-mass remains positive.

Apply the terminal local concentration decomposition to a terminal extraction of this
tail-normalized candidate.  Local vanishing, exterior escape, diffuse exterior
concentration, and loss of local compactness cannot be the sole locations of the
persistent compact local concentration in the generic residual, by
\Cref{thm:terminal-nonactive-discharge}.  Therefore the retained concentration
is located either in a finite separated family of retained concentration
profiles or in a single retained concentration terminal profile.

The finite separated-family case contradicts
\Cref{thm:no-finite-separated-profile-family}.  In the single-profile case,
\Cref{prop:no-multi-implies-atomic} makes the terminal profile terminally
indecomposable.  Then \Cref{lem:no-atomic-active} excludes it.  Both possible
locations of the retained compact velocity mass are impossible.  Hence no
generic residual profile exists.
\end{proof}

\section{Final assembly}
\label{sec:final-assembly}

The class closures above form the following dependency chain.  The
rotational relative-equilibrium arrow is the localized Coriolis-flux argument; the
terminal local concentration, no finite separated-family, and
indecomposability arrows
are proved in this paper from local estimates:
\[
\begin{aligned}
\Cax
   &\xrightarrow{\;\Gamma\hbox{-Liouville}\;}
   \varnothing,\\
\Crot
   &\xrightarrow{\;\hbox{local flux}\;}
   \varnothing,\\
\Cstat
   &\xrightarrow{\;\hbox{terminal concentration}\;}
   \varnothing,\\
\Cgen
   &\xrightarrow{\;\hbox{tail normalization, finite separation, sequence-}L^3\;}
   \varnothing .
\end{aligned}
\]
The positive local velocity mass inherited by every normalized Seregin limit is the
common contradiction target in all four classes.

\begin{proof}[Proof of \Cref{thm:refined-residual-closure}, revisited]
The detailed reductions have now been given in the four class sections.  The
axisymmetric bounded-circulation class is empty unconditionally by
\Cref{thm:R1-exclusion}.  The rotational relative-equilibrium class is empty
by the local Coriolis-flux and recentered concentration analysis in
\Cref{thm:R2-exclusion-safe}.  The degenerate stationary-hull class with
retained local concentration is empty by the terminal local concentration theorem, through
\Cref{thm:R3S-R3-empty}.  The generic terminal residual is empty by the
finite separated-family and indecomposability argument in
\Cref{thm:R4-final-closure}.
Since the refined decomposition exhausts the residual class by definition, these four
local exclusions imply that the residual class is empty.
\end{proof}

\begin{remark}[What is imported and what is proved here]
\label{rem:imported-vs-proved}
The paper imports from Paper 0 only local positive concentration and the local
Type I/Type II dichotomy, imports from Paper I only the Type I Seregin
extraction and residual criterion, imports from Paper II only the results proved
there, and imports from Paper III only the closed
structured alternatives.  The terminal finite separated-family exclusion is not imported:
it is implemented above by local concentration analysis.  The rotational
relative-equilibrium class is not closed by a global Coriolis--Leray Liouville
theorem for bounded profiles; it is closed by the localized flux identity
and the same recentered local concentration analysis.
\end{remark}

\begin{thebibliography}{99}

\bibitem{DuranPaper0}
G.~Duran Ballester,
\emph{Local CKN concentration and the Type I/Type II dichotomy for
Navier--Stokes},
preprint, 2026.

\bibitem{DuranPaperI}
G.~Duran Ballester,
\emph{Type I blow-up limits and residual closure for the
three-dimensional Navier--Stokes equations},
preprint, 2026.

\bibitem{DuranPaperII}
G.~Duran Ballester,
\emph{Uniformly tight ancient Navier--Stokes dynamics and endpoint \(L^3\)
rigidity},
preprint, 2026.

\bibitem{DuranPaperIII}
G.~Duran Ballester,
\emph{Structured ancient solution exclusions for Type I Navier--Stokes blow-up
limits},
preprint, 2026.

\bibitem{AlbrittonBarker2019}
D.~Albritton and T.~Barker,
\emph{On local Type I singularities of the Navier--Stokes equations and
Liouville theorems},
J. Math. Fluid Mech. 21 (2019), Article 43.

\bibitem{CKN1982}
L.~Caffarelli, R.~Kohn, and L.~Nirenberg,
\emph{Partial regularity of suitable weak solutions of the Navier--Stokes
equations},
Comm. Pure Appl. Math. 35 (1982), 771--831.

\bibitem{ESS2003}
L.~Escauriaza, G.~Seregin, and V.~Sverak,
\emph{\(L_{3,\infty}\)-solutions of Navier--Stokes equations and backward
uniqueness},
Russian Math. Surveys 58 (2003), 211--250.

\bibitem{KNSS2009}
G.~Koch, N.~Nadirashvili, G.~Seregin, and V.~Sverak,
\emph{Liouville theorems for the Navier--Stokes equations and applications},
Acta Math. 203 (2009), 83--105.

\bibitem{LRZ2022}
Z.~Lei, J.~Ren, and X.~Zhang,
\emph{A Liouville theorem for the axially symmetric Navier--Stokes equations},
J. Differential Equations 313 (2022), 213--238.

\bibitem{LZZ2017}
Z.~Lei, Q.~Zhang, and N.~Zhao,
\emph{Improved Liouville theorems for axially symmetric Navier--Stokes
equations},
J. Math. Fluid Mech. 19 (2017), 273--285.

\bibitem{NRS1996}
J.~Necas, M.~Ruzicka, and V.~Sverak,
\emph{On Leray's self-similar solutions of the Navier--Stokes equations},
Acta Math. 176 (1996), 283--294.

\bibitem{Seregin2012}
G.~Seregin,
\emph{A certain necessary condition of potential blow up for Navier--Stokes
equations},
Comm. Math. Phys. 312 (2012), 833--845.

\end{thebibliography}

\end{document}
